% Modernised reading — fonds Alexandre Grothendieck, shelfmark 67
% (pages 1-100, the whole folder).
%
% This is an INTERPRETATION, not a transcription. It re-reads the pages in
% today's notation and vocabulary, states the hypotheses the manuscript leaves
% implicit, and drops the critical apparatus so that the mathematics reads
% straight through. Everything the brackets and underlines carried moves into a
% footnote. It is held to being mathematically correct as it stands; where the
% manuscript is loose or mistaken, the true statement is what appears here and
% the footnote says what the page has. In any dispute of substance the
% transcription governs: whoever cites Grothendieck must cite batch-01.fr.tex
% (pp. 1-20), batch-02.fr.tex (pp. 21-40), batch-03.fr.tex (pp. 41-60),
% batch-04.fr.tex (pp. 61-80) or batch-05.fr.tex (pp. 81-100).
%
% Pass: Opus 5.5 (claude-opus-5-5), 2026-09-22 — first pass, unchecked against the pages by a human.
% The folder was read whole, all five batches together; this is one document
% for the shelfmark. It was made on Opus 5.5 and has not been measured against
% the reference reading of folder 115 (made on Opus 5).
% Revised 2026-09-23 (Opus 5.5): the order of the canonical-metrics leaves
% restated on the circled numbers as re-read on the facsimile (pp. 41-43 are
% his 16-18, pp. 44-45 his inserts 1′ and 2′), in the header and in « L'ordre
% des feuillets » and « La numérotation »; the second series' numbering
% (pp. 53-81) stated leaf by leaf; the Tate-curve footnote of section II no
% longer says Roquette published Tate's manuscript (Roquette 1970 expounds
% the theory, Tate's text appeared in 1995; cited from memory); the p. 100
% footnote no longer says the page writes « à » above the struck « sans ».
%
% Page numbers. Every page reference is the transcription's \page{}: the
% leaf's rank in the folder, confirmed by the Montpellier footer. From p. 22
% on, the archivists' pencil numbers run three ahead (p. 22 is pencilled 25,
% p. 100 is pencilled 103). His own sheet numbers are cited as « son feuillet
% n ».
%
% Order of the leaves (established here from the text, not from the numbers).
%  - pp. 26-43 are his sheets 1-18, circled numbers top left (re-read on the
%    facsimile 2026-09-23: pp. 41-43 carry 16, 17, 18). P. 41 continues p. 40,
%    whose last sentence (« la deuxième ») runs on into it.
%  - pp. 44 and 45, on another paper, are inserts, circled « 1′ » (probable;
%    « 17 » not wholly excluded) and « 2′ ». Where they go is inferred from
%    the text: p. 44 between p. 26 and p. 27 — it carries Corollaire 3 after
%    the Cor. 1-2 of p. 26 and ends « On trouve donc de plus », which p. 27
%    continues « les surfaces suivantes » (p. 27 is noted as lacking its
%    lead-in); p. 45 between p. 27 and p. 28 — it ends « La constante
%    multiplicative », and p. 28 opens « déterminée par la condition que
%    l'aire totale de X soit = 1 ». The primes agree with those places.
%    Reading order: 26, 44, 27, 45, 28-43.
%  - pp. 53-81 are one run, a second series numbered at the head of the
%    rectos: 1-4 on pp. 53, 55, 57, 59, the digit on p. 60 scratched out,
%    5-14 on pp. 61-79, 15 on p. 81. Dated 11.1.84 (p. 53), 12.1. (margin
%    p. 60 and p. 69), 13.1. (p. 73). P. 81 finishes p. 80's sentence.
%
% Spine. The folder is after one thing, which the letter to Bers (pp. 15-20)
% states: how the Teichmüller space of curves with nu marked points sits in
% the Teichmüller space of bordered surfaces (Thurston's, with boundary
% lengths), through the surgery that caps each hole with a canonical disc.
% Stations: the 1977-78 programme of surgeries (pp. 1-6); the Nagoya
% envelopes, an elliptic curve with a closed geodesic (pp. 8-13, 1977); the
% letter (pp. 15-21, 15.4.1984); canonical metrics, which make the capping
% canonical (pp. 23-45, read in the order above); Teichmüller isomorphisms
% and anabelian questions (pp. 47-50, « 15.1. »); the January 1984 run:
% foundations of cutting and gluing (pp. 53-60), discs, shadows and
% admissible systems with a Conjecture (pp. 60-69), Lemme 4 (pp. 69-73),
% passage to stable curves (pp. 73-81); M_{0,4} and the octahedron
% (pp. 83-88); real elliptic curves (pp. 90-100, « Déc 83 »). The links
% between stations that are not written on the pages are ours and are said so.
%
% Notation fixed for the document.
%  - « surface conforme »: possibly non-orientable, possibly with boundary
%    (a Klein surface); « surface holomorphe »: a Riemann surface.
%  - \mathbb{H} open upper half-plane, \overline{\mathbb{H}} closed (the page
%    writes H-bar, then H from p. 55); \mathbb{D} closed unit disc,
%    \mathring{\mathbb{D}} open; \mathbb{U} the circle group.
%  - X^{or}: the complex double (orientation double, folded along the
%    boundary), with anti-involution sigma; p. 41 writes \mathfrak{X}.
%  - T_{g,nu}: Teichmüller space of curves of type (g,nu) (the letter's
%    \widetilde{M}_{g,nu}); TB_{g,nu}: bordered, lengths > 0 (the letter's
%    \widetilde{MB}^{o}, p. 66's TB); \overline{TB}_{g,nu}: generalised
%    boundary, lengths >= 0 (the letter's \widetilde{MB}, same symbol as the
%    former on its p. 15).
%  - Teichmüller groups (pp. 49-50): \Gamma_{g,nu} with his decorations
%    + (orientation-preserving), - (the reversing coset), ! (pure); the page
%    writes T. Groupoids keep \mathcal{T}.
%  - o(D): the shadow (« ombre ») of a disc, in T_{X,a}; the letter's
%    \Delta_i. Marked points a_i (the letter's s_i).
%  - « structure torique » on a circle: a simply transitive action of U.
%  - pp. 84-88: the page's underlined antipodism a and torsor t are written
%    \mathsf{a}, \mathsf{t} (an underline would read as the apparatus).
%  - « admissible » means two things: an equivalence relation (1977, p. 2) and
%    a system of discs (1984). Kept apart in the text.
%  Collisions not corrected: T is also a torus (pp. 31-36) and a form of G_m
%  (p. 12); V is a complex plane (p. 31), the Klein group (pp. 83-88) and a
%  complex line (p. 100); Sigma is a sphere (letter), a line (p. 87) and finite
%  sets (p. 79).
%
% Substantive departures, with pages.
%  p. 6    3/2 (lambda-1) -> 1/2 (lambda-1) (his next line agrees); disc faces
%          contribute +3/2 each to N-(3g-3), dropped on the page; N is a real
%          count, 3g-3 a complex dimension (real comparison given); « g(f)=0
%          toujours » for g=1 needs a hypothesis (supplied).
%  p. 10-12 (6) « formes de Gl(1) » read as groups isomorphic to PGL_2 = Aut P.
%  p. 8-12 closed geodesic: simple (supplied, so that a is primitive).
%  p. 15   rigidification by « homotopy equivalence »: needs orientation and
%          the peripheral structure (supplied).
%  p. 18   « usual euclidean metric » with 0,1,oo equidistant: it is the unique
%          spherical metric invariant under the anharmonic group (by p. 36).
%  p. 19   the pants model is FALSE when two discs are degenerate (extremal
%          length; the argument is ours).
%  p. 27   « rev. universels » -> orientation double covers (P^1, C*).
%  p. 29   case 2'') concerns the affine line C, not an elliptic curve.
%  p. 33   his margin « c'est l'inverse ! » adopted; p. 34 « normalisent » ->
%          « centralisent ».
%  p. 44   X \ X' -> X' \ X.
%  p. 53   Th. 1: given the topology of X (supplied). p. 54: « lemme de
%          Schwarz » = reflection principle; holomorphy inside supplied.
%  p. 55   Th. 2: morphisms sending boundary to boundary (supplied).
%  p. 55-58 Th. 3 needs no rectifiability (Carathéodory); Th. 4 with
%          « rectifiable » on both sides is not correct as it stands: the
%          matching classes are quasiarcs / quasisymmetric maps (conformal
%          welding); rectifiability gives uniqueness (removability).
%  p. 60   margin c): « continue » -> holomorphe; true for rectifiable C,
%          false for arbitrary Jordan curves.
%  p. 64   canonical toric structure on the boundary fails for (g,nu)=(0,1).
%  p. 67   the Conjecture (shadow map bijective) cannot hold: Koebe bounds the
%          shadows (ours); the letter of 15.4 replaces it by r_i < rho_i.
%  p. 70   rectifiability of C, I not needed in Lemme 4 A.
%  p. 72   f(0) = 0 -> f(1) = 0; B's converse and the Scholie a)<->b) are the
%          welding problem, not established on the page.
%  p. 84   the stabiliser of a face is C_3; the splitting is the stabiliser of
%          a pair of antipodal faces.
%  p. 85   field: quadratically closed, char != 2 (« réelle » cannot be).
%  p. 91   sigma(x,y) = x - y -> (x,-y).
%  p. 97   a = sqrt 3 -> a = 3 (E'_3, E'_{1/3}), as tau = (1+i sqrt 3)/2 needs.
%  p. 100  non-trivial torsors give Klein bottles WITHOUT boundary.
%  p. 50   T^+_{0,1} -> Gamma^+_{1,1}. p. 49: r_i^2 = 1 implicit.
%
% Announced and never established on the pages: the mixed operation and 6)
% of p. 3; the moduli count of pp. 4-6 (its normalisation is illegible);
% Th. 4 and its Lemme 3 (« on doit avoir »); the fibration TB -> T (p. 66);
% the Conjecture p. 67; the Question p. 77 and the (0,4) key case; the
% étaleness of p. 81; « À examiner » a)-c) p. 48; the three presentations
% p. 49; the octahedral X(I) of p. 86 (illegible); p. 94's presumption
% (refuted by p. 97 itself).
\documentclass[11pt,a4paper]{article}
% Le préambule commun à toutes les transcriptions du fonds Grothendieck.
%
% Il tient en une idée : **l'incertitude se compose, elle ne se résout pas.**
% Une transcription qui lisse un mot douteux a détruit la seule chose pour
% laquelle on la faisait — un lecteur doit pouvoir savoir, à chaque endroit,
% ce qui était sur le feuillet et ce qui a été ajouté. Les six macros qui
% suivent sont donc l'appareil critique, et non de la décoration.
%
% Les mêmes commandes sont reconnues par `scripts/render.mjs`, qui produit la
% vue de lecture du site. Une macro ajoutée ici doit l'être aussi là-bas,
% faute de quoi le rendu échouera bruyamment — ce qui est voulu : un rendu
% silencieusement faux est pire qu'un rendu absent.

\ProvidesPackage{grothendieck}[2026/08/08 Transcriptions du fonds Grothendieck]

\usepackage{amsmath,amssymb,amsthm}
\usepackage{mathrsfs}
\usepackage{stmaryrd}
% Les diagrammes commutatifs sont partout dans ce fonds : c'est la langue
% même de la géométrie algébrique de Grothendieck, et une transcription qui
% les rendrait en prose ne transcrirait rien.
\usepackage{tikz-cd}
% Français par défaut — c'est la langue des feuillets — mais l'anglais est
% chargé aussi : la traduction et le résumé sont des documents anglais, et la
% typographie française (espaces avant les deux-points) y serait une faute.
% Ces documents basculent par \selectlanguage{english} en tête de corps.
\usepackage[english,french]{babel}
\usepackage{fontspec}
\usepackage{geometry}
\geometry{a4paper,margin=25mm,marginparwidth=28mm}
\usepackage{marginnote}
\usepackage{xcolor}
% ulem, not soul. `soul`'s \st and \ul are built on a character-by-character
% scan that XeLaTeX's font machinery breaks: the document fails with an
% undefined control sequence rather than degrading. ulem does the same two jobs
% and is Unicode-engine safe. `normalem` stops it hijacking \emph, which the
% transcriptions use for Grothendieck's own emphasis.
\usepackage[normalem]{ulem}
\usepackage{hyperref}
\hypersetup{colorlinks=true,linkcolor=black,urlcolor=black,pdfborder={0 0 0}}

% --- Métadonnées du lot -----------------------------------------------------
% Reprises telles quelles par le script de rendu, qui les affiche en tête de
% la vue de lecture. Elles fixent ce que le fichier prétend couvrir : sans
% elles, une transcription flotte sans point d'attache dans un fonds de seize
% mille feuillets.

\newcommand{\folder}[1]{\gdef\grothFolder{#1}}
\newcommand{\batch}[1]{\gdef\grothBatch{#1}}
\newcommand{\leaves}[2]{\gdef\grothFirst{#1}\gdef\grothLast{#2}}
\newcommand{\foldertitle}[1]{\gdef\grothFoldertitle{#1}}
\gdef\grothFolder{?}\gdef\grothBatch{?}\gdef\grothFirst{?}\gdef\grothLast{?}\gdef\grothFoldertitle{}

% --- Le résumé --------------------------------------------------------------
%
% Il ouvre la lecture modernisée, sous le titre « Résumé », et il est composé
% en petit. Ce n'est pas un ornement : c'est le seul passage du document qui
% ne soit pas des mathématiques, et un lecteur doit pouvoir le reconnaître
% avant de l'avoir lu — puis le sauter s'il connaît déjà le sujet, ou s'y
% arrêter s'il ne le connaît pas. Le filet qui le suit marque la reprise du
% corps.

\newenvironment{resume}
  {\small\setlength{\parskip}{0.45em}}
  {\par\medskip\hrule\medskip}

% --- Mots-clés ---------------------------------------------------------------
%
% En anglais, en clôture du résumé de la lecture modernisée : le vocabulaire
% moderne sous lequel chercher ce que les feuillets construisent. Le manifeste
% du site (`npm run manifest`) les extrait pour étiqueter la cote entière —
% cette ligne est la source unique des tags, il n'y a pas de fichier à côté.
\newcommand{\keywords}[1]{\par\smallskip\noindent
  {\footnotesize\textsc{Keywords} --- \textit{#1}}\par}

% --- L'appareil critique ----------------------------------------------------

% Le numéro de feuillet, en marge. C'est l'ancre qui permet de régler un
% désaccord sur une formule en regardant *un* feuillet plutôt qu'une chemise
% de six cents. Le site s'en sert aussi pour tourner les pages du fac-similé
% quand on fait défiler la transcription.
\newcommand{\leaf}[1]{%
  \marginnote{\footnotesize\color{gray!70}\textsf{#1}}%
  \label{leaf:#1}\ignorespaces}

% Illisible : on ne devine pas. Un mot inventé qui se lit comme les autres est
% le pire résultat possible.
\newcommand{\ill}{\textcolor{red!60!black}{[\,\ldots\,]}}

% Lecture douteuse : le mot est donné, et signalé comme incertain.
\newcommand{\uncertain}[1]{\uline{#1}}

% Ajout de l'éditeur — une lettre manquante, un mot sous-entendu.
\newcommand{\add}[1]{\textcolor{blue!50!black}{[#1]}}

% Biffé par Grothendieck lui-même. Ce qu'il a rayé fait partie du document :
% une rature dit souvent où la pensée a changé de direction.
\newcommand{\struck}[1]{\sout{#1}}

% Note du transcripteur, sur l'état matériel du feuillet.
\newcommand{\note}[1]{\par\smallskip{\small\color{gray!60!black}\textit{[#1]}}\par\smallskip}

% Note marginale de Grothendieck, distincte du corps du texte.
\newcommand{\marginal}[1]{\par\smallskip\noindent\hspace{1em}%
  {\small\color{gray!60!black}#1}\par\smallskip}

% --- Titre ------------------------------------------------------------------

\AtBeginDocument{%
  \begin{center}
    {\large\bfseries Fonds Alexandre Grothendieck --- cote n\textsuperscript{o}~\grothFolder}\\[2pt]
    {\itshape\grothFoldertitle}\\[4pt]
    {\small feuillets \grothFirst--\grothLast\ (lot \grothBatch)}\\[2pt]
    {\footnotesize\color{gray!60!black}Transcription établie d'après le fac-similé numérisé
     par l'Université de Montpellier.\\
     Les lectures incertaines sont soulignées, les illisibles notées [\,\ldots\,],
     les ajouts entre crochets.}
  \end{center}
  \vspace{1em}}

\endinput


\folder{67}
\pages{1}{100}
\dating{[à partir de 1977]-1984}
\watermark{Édition de démonstration\\interprétation personnelle de l'œuvre}
\foldertitle{Chirurgie des surfaces conformes : notes manuscrites (s.d., 1983-1984), lettres (1984) — lecture modernisée du dossier entier}

\begin{document}

\section*{Résumé}

\begin{resume}
Ce dossier tourne autour de gestes que l'on connaît du papier et du tissu :
découper, recoudre, rapiécer. L'étoffe est ici une surface de Riemann ---
une surface sur laquelle on sait mesurer les angles, mais pas encore les
longueurs --- et la question est de savoir ce que ces gestes font à sa forme.
Grothendieck appelle cela la \emph{chirurgie des surfaces conformes}. Le mot
est sur une chemise de 1977 ou 1978, et l'essentiel des feuillets date de
décembre 1983 à avril 1984.

Le problème vient de ce qu'il existe, au début des années 1980, deux façons
de ranger les surfaces de Riemann d'un type donné. La première, algébrique,
considère une courbe fermée sur laquelle on a marqué quelques points. La
seconde, celle de Thurston, considère une surface à bord dont chaque bord est
un cercle, avec sa longueur. Chacune donne un espace de paramètres, un
\emph{espace de Teichmüller}. Grothendieck veut savoir exactement comment le
premier s'emboîte dans le second : il devrait en être le bord, là où les
longueurs des cercles tendent vers zéro et où chaque trou se referme en un
point. Il l'écrit en avril 1984 à Lipman Bers, qu'il tient pour le spécialiste
de la question, et la lettre est dans le dossier.

Le geste qui fait passer d'un espace à l'autre est de boucher chaque trou avec
un disque. Pour coudre une pièce ronde au bord d'un trou, il faut que la pièce
et le bord « tournent » de la même façon : un cercle nu n'a pas de graduation,
et il y a mille manières de l'appliquer sur le bord d'un disque. La graduation
viendra d'une \emph{métrique canonique}. Toute surface conforme, sauf la
sphère et le plan, porte une façon privilégiée de mesurer les longueurs, à
courbure constante, et une trentaine de feuillets établissent la liste
complète des cas, y compris pour les surfaces non orientables et à bord. Le
bord devient alors un cercle gradué, et l'on peut y coudre un disque gradué de
la même manière. Reste à savoir quels disques naissent ainsi --- il les appelle
\emph{admissibles} --- et comment les mesurer. Chaque disque autour d'un point
laisse dans le plan tangent en ce point une « ombre », un disque rond dont le
rayon est ce qu'on appelle aujourd'hui le rayon conforme. En janvier 1984, il
conjecture que l'ombre détermine le disque, et que toutes les ombres sont
possibles. En avril, dans la lettre, il a vu que ce n'est pas tout à fait
cela, et il propose pour la sphère percée de trois trous un modèle en
quartiers d'orange, en demandant à Bers s'il est vraisemblable ou
franchement faux. Sous la forme énoncée, ni la conjecture de janvier ni le
modèle d'avril ne peuvent tenir ; on le montre en note, par des arguments qui
sont les nôtres.

Autour de ce fil, le dossier pose les fondations et regarde plus loin. Les
fondations : découper une surface le long d'un graphe donne une surface à
bord, et recoller des bords redonne une surface. La première opération est
aujourd'hui un théorème classique de représentation conforme ; la seconde est
le problème de la \emph{soudure conforme}, dont on sait maintenant que la
bonne hypothèse n'est pas celle que proposent les pages. Plus loin : si l'on
pince jusqu'au bout toutes les coutures d'une surface, on obtient une
\emph{courbe stable} au sens de Deligne et Mumford, et l'on veut garder la
trace de la façon dont on l'a pincée. De là, les groupes fondamentaux et les
groupes de Teichmüller, et les questions de géométrie anabélienne qu'il posait
alors, contemporaines de l'\emph{Esquisse d'un programme}. Deux études
complètent le dossier : les quatre points d'une droite projective et
l'octaèdre que dessine leur groupe de symétries, qui décrivent l'espace des
modules $M_{0,4}$ ; et les courbes elliptiques réelles, c'est-à-dire les tores
munis d'un miroir, qui achèvent la classification des surfaces de genre un.

Ces pages montrent aussi une manière de travailler. Une conjecture y est
écrite, étiquetée, et aussitôt mise en doute --- « c'est probablement tout à
fait faux ! », « sauf erreur », « à vérifier » dans les marges. Dans la
lettre, Grothendieck avoue n'avoir pas regardé la situation de près, et il
soumet au spécialiste un modèle précis plutôt qu'une intuition vague :
quelque chose qu'on puisse réfuter.

Les noms sous lesquels chercher la suite sont ceux de surface de Klein,
principe de réflexion, soudure conforme, rayon conforme, coordonnées de
Fenchel--Nielsen, plombage, courbe stable, géométrie anabélienne, tour de
Teichmüller.

\keywords{conformal surgery, bordered Riemann surface, Klein surface, Schottky double, Schwarz reflection principle, Carathéodory extension theorem, conformal welding, removable set, uniformization theorem, hyperbolic metric, spherical metric, Tate curve, loxodromic transformation, Teichmüller space, pants decomposition, Fenchel–Nielsen coordinates, conformal radius, plumbing, stable curve, Deligne–Mumford compactification, anabelian geometry, mapping class group, Teichmüller groupoid, moduli space of four points on the line, octahedral group, Klein four-group, real elliptic curve, Galois cohomology}
\end{resume}

\pagerange{1}{100}
\section*{Le fil du dossier, l'ordre des feuillets, les conventions}

\emph{La numérotation.} Les renvois de page de cette lecture sont ceux de la
transcription : le rang du feuillet dans le dossier, que confirme le pied de
page imprimé par Montpellier. À partir de la page 22, les numéros portés au
crayon par les archivistes ont trois unités d'avance (la page 22 porte 25, la
page 100 porte 103) ; on ne les utilise pas. Grothendieck numérote lui-même
deux suites de feuillets : la première de 1 à 18 (pages 26 à 43), avec deux
feuillets intercalaires numérotés $1'$ et $2'$ (pages 44 et 45) ; la seconde de
1 à 15, en tête des rectos des pages 53 à 81, le numéro de la page 60 étant
raturé. On les cite comme « son feuillet $n$ ».

\emph{Les pièces et leur date.} Le dossier réunit des feuillets de trois
époques.

\begin{enumerate}
  \item \emph{Pages 1 à 13}, fin 1977 ou 1978. La chemise (page 1) porte le
  titre et « (1978 ou 77 ?) » de sa main. Un programme de chirurgie occupe les
  pages 2 à 6 (section I). Suivent trois brouillons d'une même liste, écrits
  sur des enveloppes (pages 8 à 13, section II) ; l'une d'elles porte le
  cachet de l'université de Nagoya du 20 décembre 1977, date observée sur le
  papier et non déduite.
  \item \emph{Pages 15 à 21}, la lettre dactylographiée à Lipman Bers, datée
  des Aumettes le 15 avril 1984, et le double carbone de son dernier feuillet
  (section III).
  \item \emph{Pages 23 à 45}, sans date : les métriques canoniques
  (section IV).
  \item \emph{Pages 47 à 50}, datées « 15.1. » sans année, sur du papier de
  listing imprimé en août 1982 : isomorphismes de Teichmüller et questions
  anabéliennes (section V). La page 52 est un brouillon (section VI).
  \item \emph{Pages 53 à 81}, une suite datée « 11.1.84 », puis « 12.1. » et
  « 13.1. » : les fondements de la chirurgie, les disques admissibles, le
  passage aux courbes stables (sections VII à X). Le « 15.1. » des pages 47
  à 50 appartient vraisemblablement à la même semaine de janvier 1984.
  \item \emph{Pages 83 à 88}, sans date : $M_{0,4}$ et l'octaèdre
  (section XI).
  \item \emph{Pages 90 à 100}, datées « Déc 83 », avec en marge de la
  dernière un « 9.2. » : les courbes elliptiques réelles (section XII).
\end{enumerate}

Les pages 7, 11, 14, 22, 46, 51, 82 et 89 ne portent rien pour la
mathématique (versos sans rapport, faces d'enveloppe, blancs).

\emph{L'ordre des feuillets sur les métriques canoniques.} Les pages 26 à 43
sont ses feuillets 1 à 18, numérotés et cerclés en haut à gauche ; la
dernière phrase de la page 40, son feuillet 15, se poursuit en haut de la
page 41, son feuillet 16. Les pages 44 et 45, sur un autre papier, portent en
haut des numéros cerclés qu'on lit « $1'$ » et « $2'$ » --- le second nettement,
le premier sous réserve : ce sont des feuillets intercalaires, et c'est leur
texte qui dit où ils s'insèrent. La page 44 s'insère entre les pages 26 et 27 : elle porte un
« Corollaire 3 » qui fait suite aux corollaires 1 et 2 de la page 26, et
s'achève sur « On trouve donc de plus », que la page 27, dépourvue
d'introduction, continue par « les surfaces suivantes ». La page 45 s'insère
entre les pages 27 et 28 : elle s'achève sur « La constante multiplicative »,
et la page 28, qui commence au milieu du cas 2, continue par « déterminée par
la condition que l'aire totale de $X$ soit $= 1$ ». On lit donc : 26, 44,
27, 45, 28 à 43. Ces deux raccords sont les nôtres et reposent sur le texte ;
les numéros primés, qui placent les intercalaires après ses feuillets 1 et 2,
s'accordent avec eux.

\emph{Le fil.} La lettre à Bers dit en toutes lettres ce que le dossier
cherche. Il s'agit de comprendre de façon intrinsèque comment l'espace de
Teichmüller des courbes de type $(g,\nu)$ --- genre $g$, $\nu$ points marqués
--- se place dans celui des surfaces à bord de Thurston, où chaque trou est un
cercle muni de sa longueur. L'outil est une chirurgie : boucher chaque trou
par un disque canonique.

\begin{itemize}
  \item Le programme de 1977 (section I) énumère déjà les opérations :
  découper, recoller, rétrécir ou élargir un collier, et il pose la question
  des modules d'une surface munie d'un graphe.
  \item Les métriques canoniques (section IV) rendent le bouchage canonique :
  elles graduent chaque cercle de bord.
  \item Les fondations (sections VII et IX) disent quand découper et recoller
  a un sens.
  \item Les disques admissibles et leurs ombres (section VIII), puis la
  lettre (section III), formulent la conjecture et le modèle.
  \item Le passage aux courbes stables (section X) pousse la chirurgie
  jusqu'à la dégénérescence et la relie aux groupes de Teichmüller, dont les
  pages 47 à 50 (section V) posent les questions anabéliennes.
  \item Deux études latérales complètent la classification : l'espace
  $M_{0,4}$ (section XI), pièce de base des tours de Teichmüller, et les
  courbes elliptiques réelles (section XII), qui donnent les surfaces
  conformes de genre un non orientables ou à bord.
  \item Les enveloppes de 1977 (section II) traitent le cas le plus simple
  d'un pincement : une courbe elliptique munie d'une géodésique fermée.
\end{itemize}

Les liens entre ces stations, sauf ceux que les pages écrivent (le renvoi de
la page 69 aux lemmes 2 et 3 et au théorème 4, par exemple), sont les nôtres.

\emph{Conventions.}

\begin{itemize}
  \item Une \emph{surface conforme} peut être non orientable et avoir un
  bord : c'est ce qu'on appelle aujourd'hui une surface de Klein. Une
  \emph{surface holomorphe} est une surface de Riemann.
  \item $\mathbb{H}$ est le demi-plan supérieur ouvert, $\overline{\mathbb{H}}$
  le demi-plan fermé ; $\mathbb{D}$ est le disque unité fermé,
  $\mathring{\mathbb{D}}$ le disque ouvert ; $\mathbb{U}$ est le groupe des
  complexes de module $1$.
  \item $X^{\mathrm{or}}$ est le \emph{double complexe} d'une surface conforme
  $X$, muni de son anti-involution $\sigma$. Pour une surface sans bord,
  c'est le revêtement des orientations ; pour une surface à bord, c'est le
  double de Schottky.
  \item $\mathrm{T}_{g,\nu}$ est l'espace de Teichmüller des courbes de type
  $(g,\nu)$ ; $\mathrm{TB}_{g,\nu}$ celui des surfaces à bord, longueurs du
  bord strictement positives ; $\overline{\mathrm{TB}}_{g,\nu}$ celui des
  surfaces à bord généralisé, longueurs positives ou nulles.
  \item Les groupes de Teichmüller (groupes modulaires des surfaces) sont
  notés $\Gamma_{g,\nu}$, avec ses décorations : $+$ pour la partie qui
  conserve l'orientation, $-$ pour l'autre classe, $!$ pour les éléments qui
  fixent chaque point marqué. Les pages écrivent $T$, lettre qui désigne
  ailleurs un tore ou un espace tangent.
  \item L'\emph{ombre} d'un disque $D$ centré en $a$ est un disque rond $o(D)$
  de l'espace tangent $T_{X,a}$. Une \emph{structure torique} sur un cercle
  est une action simplement transitive de $\mathbb{U}$.
  \item Le mot \emph{admissible} a deux sens, qu'on tient séparés : en 1977,
  une relation d'équivalence de recollement ; en 1984, un système de disques.
\end{itemize}

Restent des collisions de lettres, qu'on ne corrige pas : $T$ est un tore aux
pages 31 à 36 et une forme de $\mathbb{G}_m$ à la page 12 ; $V$ est un plan
complexe à la page 31, le groupe de Klein aux pages 83 à 88, une droite
complexe à la page 100.

\emph{Ce que le dossier annonce sans l'établir.} Le décompte des modules des
pages 4 à 6, dont la normalisation est illisible. Le théorème 4 et son
lemme 3, introduits par « on doit avoir ». La fibration de la page 66. La
conjecture de la page 67 et le modèle de la lettre, qui ne tiennent pas sous
la forme énoncée. La question de la page 77 et le cas-clef $(0,4)$. Le
caractère étale de la page 81. Les questions a) à c) de la page 48 et les
trois présentations de la page 49. La construction de l'octaèdre $X(I)$ à
la page 86, illisible.

\pagerange{1}{6}
\section*{I. Le programme de 1977 (pages 1 à 6)}

\pagerange{2}{3}
\subsection*{Les opérations (pages 2 et 3)}

La chemise porte le titre du dossier, « Chirurgie des surfaces conformes »,
suivi d'un mot illisible et de « (1978 ou 77 ?) ». La page 2 reprend le titre
et énumère des opérations sur une surface conforme $X$, éventuellement à
bord.

\begin{enumerate}
  \item \emph{Découpage.} On découpe $X$ le long d'un sous-complexe $K$ de
  dimension $1$, compact et rectifiable. On obtient une surface conforme à
  bord $\widetilde{X}$, dont une partie compacte du bord, $\partial'
  \widetilde{X}$, est marquée : c'est la trace des deux rives de $K$.
  \item \emph{Recollement.} Inversement, d'une surface conforme à bord, d'une
  partie compacte de son bord et d'une relation d'équivalence
  « admissible » sur celle-ci, on déduit une surface conforme.\footnote{« Au
  sens rectifiable », lu sous réserve. Ce que doit être une relation
  admissible est exactement la question que reprennent les pages 57 à 60 et
  69 à 73 ; voir la section VII, où l'on dit quelle hypothèse est la bonne.}
  \item \emph{Rétrécissage.} On choisit autour de $\partial'\widetilde{X}$ un
  collier, réunion de couronnes, et on le retire : c'est la surface rétrécie
  $\widetilde{X}^{*}$. Un homéomorphisme rectifiable $\partial'
  \widetilde{X}^{*} \simeq \partial'\widetilde{X}$, compatible aux
  composantes, transporte la relation de recollement ; on en déduit une
  nouvelle surface $X^{*}$ contenant encore $K$. Le croquis en marge marque
  l'épaisseur $\varepsilon_i$ de chaque couronne, avec
  $0 \leqslant \varepsilon_i < +\infty$.\footnote{On peut penser
  $\varepsilon_i$ comme le module conforme de la couronne retirée ; c'est
  notre lecture. Une parenthèse de la page dit à quoi revient la donnée d'un
  germe de couronne autour d'une composante (« une structure … de
  longueur totale $1$ ») ; elle se lit mal et nous ne la reconstituons pas.}
  \item \emph{Élargissage.} L'opération inverse : on ajoute une couronne au
  lieu d'en retirer une.
  \item \emph{Opération mixte.} Dans un voisinage tubulaire de
  $\partial'\widetilde{X}$, on choisit des cercles $\Gamma_i$ et on retire
  ce qu'ils entourent. La marge note que c'est la composée d'un rétrécissage
  et d'un élargissage, dans l'ordre qu'on veut.\footnote{Le nom de cette
  opération, souligné, commence par un p et finit en -age ; il revient aux
  alinéas 4') et 6) et n'a pas pu être lu.}
\end{enumerate}

Un alinéa 4') affirme que ces opérations, faites le long de composantes
connexes distinctes de $K$, commutent. L'alinéa 6) considère une composante
de $K$ qui est une courbe fermée isolée, plongée de façon analytique réelle :
l'opération revient alors à retirer une couronne autour de la courbe, ou à en
insérer une, en recollant ses deux bords sur un cylindre.\footnote{L'alinéa 6)
est aux trois quarts incertain, et la phrase qui le résume ici est une
reconstitution. Ainsi comprise, l'insertion d'un anneau le long d'une courbe
fermée est ce qu'on appelle aujourd'hui une greffe (\emph{grafting}) ; le
rapprochement est le nôtre, et la page ne le fait pas.}

\pagerange{4}{6}
\subsection*{Les modules d'une surface munie d'un graphe (pages 4 à 6)}

L'alinéa 7) pose la question des modules d'un triple $S \subset K \subset X$,
où $X$ est une surface compacte orientée sans bord, $K$ un graphe tracé sur
$X$ et $S$ un ensemble fini de points de $K$, sans doute ses sommets. On fixe
un triple topologique de référence $S_0 \subset K_0 \subset X_0$, et l'on
\emph{rigidifie} $(S \subset K \subset X)$ par une classe d'isotopie
d'homéomorphismes avec la situation de référence : c'est le procédé qui
définit l'espace de Teichmüller. Les automorphismes disparaissent ainsi, sauf
quand $X_0$ a des composantes de genre $0$ ou $1$.

Pour obtenir un véritable espace de modules, il faut normaliser la forme des
arêtes. Grothendieck envisage d'abord d'exiger que les arêtes soient des arcs
de cercle pour la structure projective canonique, puis note que certaines
« cartes » ne satisfont pas cette condition ; il envisage des paramétrages
géodésiques dans le disque hyperbolique ; il conclut, en haut de la page 6,
qu'il vaut mieux exiger que les arêtes soient \emph{horocycliques par
morceaux}.\footnote{La page 4 s'achève sur « Il », et la page 6 commence par
« vaut mieux » : la page 5, un feuillet de croquis, s'intercale. Le mot qui
qualifie les « cartes » est illisible, et plusieurs mots de l'argument le
sont aussi ; la normalisation finalement retenue n'est donc connue que par son
nom.} La page 5 est une feuille de croquis. On y lit les notations du
découpage : un revêtement $\widetilde{K} \to K$ qui dédouble $K$ hors d'un
sous-complexe $L$, avec $|\widetilde{K}| \smallsetminus |\widetilde{L}|
\simeq |K| \smallsetminus |L|$, et $\widetilde{K} = \varinjlim_{x}K_{x}$ où
$K_x$ est l'ensemble des cellules sous $x$. On y voit aussi des pantalons et
une sphère découpée en fuseaux.

Vient un décompte. On suppose $X$ connexe de genre $g$, $K$ connexe, et l'on
note $F$ l'ensemble des faces, c'est-à-dire des composantes connexes de
$X \smallsetminus K$. Chaque face $f$ est une surface ouverte de genre $g(f)$
avec $\nu(f)$ bouts, et $\lambda$ est le premier nombre de Betti de
$K$.\footnote{Le mot qui définit $\lambda$ en marge est illisible ; la
formule d'Euler qui suit n'est juste que si $\lambda$ est le rang de
$H_1(K)$, et c'est ainsi qu'on le lit.} Comme la caractéristique d'Euler à
support compact d'une face vaut $2 - 2g(f) - \nu(f)$, on a
\[
2 - 2g = \sum_{f \in F}\bigl(2 - 2g(f) - \nu(f)\bigr) + (1 - \lambda),
\]
d'où
\[
3g - 3 = \sum_{f \in F}\Bigl(3g(f) - 3 + \tfrac{3}{2}\,\nu(f)\Bigr)
+ \tfrac{3}{2}(\lambda - 1).
\]
La page passe par $g - 1 = \frac12\sum_f\bigl(2g(f)-2+\nu(f)\bigr) +
\frac12(\lambda-1)$.\footnote{La page écrit $\frac32(\lambda-1)$ dans cette
ligne, le $\frac32$ récrit sur un autre chiffre ; la ligne précédente donne
$\frac12(\lambda-1)$, et la ligne suivante, qui multiplie par $3$, est juste
avec $\frac12$.}

Grothendieck compte alors les modules réels du triple, sous l'hypothèse que
$\pi_0(K_0) \to \pi_0(X_0)$ soit surjectif, et propose
\[
N = \sum_{f \text{ non disque}} \bigl(6g(f) + 4\nu(f) - 6\bigr).
\]
Pour une face de genre $g(f)$ à $\nu(f)$ bords, $6g(f) - 6 + 3\nu(f)$ est la
dimension réelle de l'espace de Teichmüller des surfaces à bord ; on peut lire
le terme supplémentaire $\nu(f)$ comme un paramètre réel par bord, lecture qui
est la nôtre. C'est la
normalisation des arêtes qui justifierait ce décompte, et elle n'est pas
lisible ; on donne donc $N$ comme le décompte qu'il propose. Pour $g
\geqslant 2$ on trouve
\[
N - (3g - 3) = \sum_{f \text{ non disque}}\Bigl(3g(f) + \tfrac52\,\nu(f) -
3\Bigr) + \tfrac32\, d + \tfrac32(1 - \lambda),
\]
où $d$ est le nombre de faces qui sont des disques. La page demande si cette
quantité est positive.\footnote{La page n'a pas le terme $\frac32 d$ : une
somme sur les faces qui sont des disques, écrite devant le dernier terme, est
biffée, et une ligne en dessous, biffée elle aussi, en parle encore. Or chaque
face disque contribue $-\bigl(3\cdot 0 - 3 + \frac32\bigr) = \frac32$ au
membre de droite. Par ailleurs, $N$ compte des modules \emph{réels}, tandis
que $3g - 3$ est la dimension \emph{complexe} de l'espace des modules des
courbes de genre $g$. La comparaison entre dimensions réelles donne
$N - (6g - 6) = \sum_{f \text{ non disque}}\nu(f) + 3d - 3(\lambda - 1)$.
On garde la question de la page telle qu'elle est posée.} Pour $g = 1$, il
note que toutes les faces sont de genre $0$ et que chaque terme $4\nu(f) - 6$
d'une face non disque vaut au moins $2$.\footnote{« $g(f) = 0$ toujours »
demande une hypothèse : si $K$ est contenu dans un disque de $X$ (un arbre,
par exemple), une face est un tore troué, de genre $1$. L'affirmation est
vraie dès qu'un cycle de $K$ n'est pas homologue à zéro dans $X$ : une face de
genre $1$ contiendrait un tore à un trou, dont le complémentaire dans $X$ est
un disque, et $K$ y serait contenu.}

\pagerange{8}{13}
\section*{II. Courbe elliptique et géodésique fermée (pages 8 à 13)}

Trois brouillons successifs d'une même liste, aux pages 8, 10 et 12, écrits
sur des enveloppes. L'une porte au recto (page 13) le cachet de Nagoya du 20
décembre 1977. La version de la page 12, la plus complète, énumère six
données qui se correspondent deux à deux ; les doubles flèches entre les
numéros sont les siennes. On les donne sous leur forme orientée, celle des
pages 8 et 10.

\begin{enumerate}
  \item Une courbe elliptique complexe $X$ et une géodésique fermée simple
  orientée $\Gamma$ de sa métrique plate, à translation près.\footnote{« Simple »
  est ajouté : une géodésique parcourue $k$ fois correspond à $k$ fois une
  classe primitive, et la correspondance avec (3) exige la classe primitive.
  La page 12 donne la version non orientée (géodésique à translation près,
  sous-groupe facteur direct de rang $1$).}
  \item $X$ et un sous-groupe fermé $\Gamma_0 \simeq \mathbb{R}/\mathbb{Z}$
  de $X$, orienté : un « tore réel ». C'est le translaté de $\Gamma$ passant
  par l'origine.
  \item $X$ et une classe primitive $a \in \pi_1(X) = \Pi$, où
  $X = V/\Pi$ ; en version non orientée, un sous-groupe facteur direct
  $Z = \mathbb{Z}a$ de rang $1$. La variante (3') dit la même chose sans
  origine : une droite complexe $V$, un réseau $\Pi$, un facteur direct
  $Z \subset \Pi$ de rang $1$, et un torseur $X$ sous $V/\Pi$.
  \item Un groupe $T$ isomorphe à $\mathbb{G}_m$ --- une « forme » de
  $\mathbb{G}_m$ ---, un sous-groupe discret $Z \subset T(\mathbb{C})$
  isomorphe à $\mathbb{Z}$, et un torseur $\Delta$ sous $T$ ; alors
  $X = \Delta/Z$. Concrètement, $T = V/Z \simeq \mathbb{C}^{*}$ par
  $z \mapsto \exp(2i\pi z/a)$, et $X \simeq \mathbb{C}^{*}/b^{\mathbb{Z}}$
  avec $0 < |b| < 1$.
  \item Une droite projective $P$ et un automorphisme $u$ de $P$ non
  parabolique --- à deux points fixes distincts --- dont les valeurs propres
  $\lambda, \lambda^{-1}$ (d'un relèvement à $\mathrm{SL}_2$) vérifient
  $|\lambda| \neq 1$. Alors $X = (P \smallsetminus \mathrm{Fix}(u))/
  u^{\mathbb{Z}}$. C'est (4) complété : $P = \Delta \cup \{0, \infty\}$, et
  $u$ est l'homothétie de rapport $b$ ou $b^{-1}$.
  \item Un groupe $G$ isomorphe à $\mathrm{PGL}_2$ et un sous-groupe discret
  $Z \subset G(\mathbb{C})$, isomorphe à $\mathbb{Z}$ et non unipotent,
  c'est-à-dire formé d'éléments semi-simples, contenu dans un tore maximal
  $T$.\footnote{La page 12 écrit « formes $G$ de $\mathrm{Gl}(1)_{\mathbb{C}}$ »,
  et le palimpseste de la page 10 parle de sous-groupe de Borel
  $T \subset B \subset G$ et de groupe déployé. Or $\mathrm{GL}_1$ n'a ni
  élément unipotent non trivial ni Borel distinct de lui-même ; tout ce qui
  est dit n'a de sens que pour $G = \mathrm{Aut}(P) \simeq \mathrm{PGL}_2$,
  ce qui est aussi ce que demande la correspondance avec (5). La lecture est
  la nôtre. Sur $\mathbb{C}$, toute forme de $\mathrm{PGL}_2$ est triviale :
  « forme » veut dire ici « groupe isomorphe, sans isomorphisme choisi ».}
  La droite $P$ se retrouve comme la variété des sous-groupes de Borel de
  $G$, et les deux points fixes de $u$ comme les deux Borels qui contiennent
  $Z$.
\end{enumerate}

Chaque donnée porte un ensemble à deux éléments, et la page écrit la chaîne
de bijections canoniques entre eux : les deux orientations de $\Gamma$, celles
de $\Gamma_0$, les deux générateurs de $Z$, les deux points fixes de
$u$.\footnote{Sous les termes de la chaîne, les numéros cerclés des données.
Le dernier terme, « $\mathrm{Inv}(T)$ », est une lecture douteuse ; on peut y
voir les deux isomorphismes $T \simeq \mathbb{G}_m$, ce qui s'accorde avec le
reste. La suite de la chaîne, en marge, se lit « $\simeq \mathrm{Fix}(u) \simeq$ », suivie d'un dernier terme incertain.} La page 10 écrit une autre
chaîne, pour $\Gamma_0$ fixé : la famille des géodésiques parallèles à
$\Gamma_0$ est le cercle $X/\Gamma_0$. Dans le modèle multiplicatif, ces
géodésiques sont les images des cercles $|z| = r$, c'est-à-dire des cercles
de $P$ qui séparent les deux points fixes de $u$, pris à l'action de $u$
près.\footnote{La page écrit $\Delta/T_{\mathbb{C}}\cdot Z$ ; le quotient qui
paramètre les cercles $|z| = r$ modulo $b^{\mathbb{Z}}$ est
$\Delta/(\mathbb{U}\cdot Z) \simeq \mathbb{R}_{>0}/|b|^{\mathbb{Z}}$. Une
note en petit ajoute que ces cercles sont ceux qui admettent
$\mathrm{Fix}(u)$ comme « bicentre ».}

Au recto de l'enveloppe, la page 13 vérifie la condition de (5) :
\[
|\lambda| \neq 1 \iff \lambda + \lambda^{-1} \notin [-2, 2],
\]
avec le calcul qu'il faut. Si $|\lambda| = 1$, alors $\bar\lambda =
\lambda^{-1}$ et $\lambda + \lambda^{-1}$ est un réel de $[-2,2]$ ;
inversement, si $\lambda + \lambda^{-1} = \alpha \in [-2,2]$, le discriminant
$\alpha^2 - 4$ est négatif ou nul et $|\lambda| = 1$. Une telle $u$ est ce
qu'on appelle \emph{loxodromique}.

Ce qu'on reconnaît aujourd'hui dans cette liste est l'uniformisation de Tate
(ou de Jacobi) $X \simeq \mathbb{C}^{*}/q^{\mathbb{Z}}$, attachée au choix
d'un cycle $a$. Pincer $\Gamma$ jusqu'à un point, c'est faire tendre $q$ vers
$0$ et obtenir une cubique nodale.\footnote{Le nom de courbe de Tate vient
d'un manuscrit de Tate de 1959, longtemps resté inédit. Roquette en a exposé
la théorie en 1970 (\emph{Analytic theory of elliptic functions over local
fields}), et le texte de Tate lui-même n'a paru qu'en 1995 (« A review of
non-archimedean elliptic functions »). Le quotient
$\mathbb{C}^{*}/q^{\mathbb{Z}}$ remonte à Jacobi. Ces références sont citées
de mémoire, et la page ne cite rien. Le lien avec les dégénérescences des pages 73 à 81 est le nôtre.}

\pagerange{15}{21}
\section*{III. La lettre à Lipman Bers (pages 15 à 21)}

La lettre est dactylographiée, en anglais, datée des Aumettes le 15 avril
1984. Les lettres grecques y sont portées à la main, et elle n'est pas signée
sur ces feuillets. Le feuillet de la page 21 est un double carbone du
dernier, celui de la page 20. On la suit dans son ordre.

\pagerange{15}{16}
\subsection*{Ce que la lettre cherche (pages 15 et 16)}

Grothendieck dit animer avec Yves Ladegaillerie, son ancien élève, un
« microséminaire » sur les espaces et groupes de Teichmüller. Ses propres
motivations viennent de la géométrie algébrique, celles de Ladegaillerie de la
topologie des surfaces.\footnote{Les travaux de Ladegaillerie portent sur les
plongements de graphes dans les surfaces et leurs classes d'isotopie, sujet
voisin des pages 24 à 30 du dossier 161-6. La lettre n'en dit pas plus que
ce qui est rapporté ici.} Il s'adresse à Bers comme au principal expert de
l'approche de l'espace de Teichmüller par la géométrie hyperbolique de
Thurston.

On fixe une surface topologique compacte orientée $X_0$, de genre $g$, avec
$\nu$ trous, et $2g - 2 + \nu > 0$. L'objet qui l'intéresse d'abord est
l'espace de Teichmüller algébrique $\mathrm{T}_{g,\nu}$. Il classe les
courbes algébriques lisses $X$ sur $\mathbb{C}$, de genre $g$, munies de
$\nu$ points distincts $S$ et d'une rigidification de Teichmüller de
$(X,S)$ : une classe d'homotopie d'homéomorphismes de $X_0$, privée de son
bord, sur $X \smallsetminus S$, qui conservent l'orientation.\footnote{La
lettre dit « une équivalence d'homotopie entre $X_0$ et
$X \smallsetminus S$ ». Pour $\nu \geqslant 1$, c'est trop large : il y a des
équivalences d'homotopie qui ne respectent pas les pointes, et pour le
pantalon ($g = 0$, $\nu = 3$), dont le groupe fondamental est libre de rang
$2$, elles forment $\mathrm{GL}_2(\mathbb{Z})$ alors que le groupe modulaire
est fini. Il faut exiger qu'elle conserve l'orientation et le système des
classes périphériques ; d'après le théorème de Dehn--Nielsen--Baer, cela
revient à la donnée qu'on écrit ici.} Il est homéomorphe à $\mathbb{C}^{d}$,
avec $d = 3g - 3 + \nu$. Le tilde de sa notation
$\widetilde{M}_{g,\nu}$ rappelle qu'il est le revêtement universel d'un objet
plus fin, qui l'intéresse davantage : la multiplicité modulaire, ou champ au
sens de Deligne et Mumford, $M_{g,\nu}$.\footnote{Au sens des orbifolds :
$\mathrm{T}_{g,\nu}$ est contractile et le champ analytique associé à
$M_{g,\nu}$ est le quotient de $\mathrm{T}_{g,\nu}$ par le groupe modulaire.
La lettre écrit « Mumford--Deligne ». Que l'espace de Teichmüller soit
homéomorphe à une boule était classique en 1984.}

Thurston considère un autre espace, celui des surfaces conformes compactes
orientées \emph{à bord} : c'est $\mathrm{TB}_{g,\nu}$, homéomorphe à
$\mathbb{C}^{d} \times (\mathbb{R}^{*}_{+})^{\nu}$. Le facteur
supplémentaire est donné par les longueurs $l_i$ des composantes du bord pour
la métrique hyperbolique canonique, celle pour laquelle le bord est
géodésique. Grothendieck veut préciser la relation entre les deux espaces.
L'idée naturelle est de voir une courbe munie de $\nu$ points comme la limite
d'une surface à bord dont toutes les longueurs $l_i$ tendent vers zéro. On
plonge donc les deux dans un troisième : l'espace
$\overline{\mathrm{TB}}_{g,\nu}$ des surfaces à \emph{bord généralisé}, où
certaines composantes du bord peuvent être réduites à un point. Il est
homéomorphe à $\mathbb{C}^{d} \times (\mathbb{R}_{+})^{\nu}$ : c'est une
variété à bord (à coins, en analytique réel), dont $\mathrm{T}_{g,\nu}$ est la
face où toutes les longueurs sont nulles.\footnote{La lettre note du même
symbole $\widetilde{MB}_{g,\nu}$ l'espace de Thurston et l'espace à bord
généralisé, et écrit ensuite $\widetilde{MB}^{o}$ pour le premier. On
distingue ici $\mathrm{TB}$ et $\overline{\mathrm{TB}}$. Les coordonnées
qu'elle appelle « standard », longueurs et torsions relatives à une
décomposition en pantalons, portent aujourd'hui le nom de Fenchel et
Nielsen.}

Ce qu'il cherche est une compréhension \emph{intrinsèque}, indépendante de
toute décomposition de $X_0$ en pantalons. Il attend une rétraction
naturelle de $\overline{\mathrm{TB}}$ sur $\mathrm{T}$, compatible au groupe
modulaire, à fibres homéomorphes à $(\mathbb{R}_{+})^{\nu}$. Plus
précisément, il attend une action du semi-groupe $(\mathbb{R}_{+})^{I}$,
indexé par les trous, telle que le sous-groupe $(\mathbb{R}^{*}_{+})^{I}$
agisse librement sur $\mathrm{TB}$, que $\mathrm{T}$ soit le quotient, et
que chaque fibre devienne isomorphe à $(\mathbb{R}_{+})^{I}$ dès qu'on y
choisit une origine. En coordonnées de pantalons, l'action évidente
multiplie chaque longueur $l_i$ par $\lambda_i$, et le multiplicateur nul
donne la rétraction ; mais rien ne dit que cette action soit intrinsèque.

\pagerange{16}{18}
\subsection*{Boucher les trous, et les disques admissibles (pages 16 à 18)}

Il y a une description géométrique directe de la rétraction, par chirurgie.
On bouche chaque trou qui est un vrai cercle en y collant le cône sur ce
cercle. Le cercle, muni de sa structure riemannienne, est un cercle gradué,
et le cône porte une structure conforme canonique. On obtient une surface
compacte $X^{\wedge}$ sans bord, munie de $\nu$ points. Mais on obtient plus :
un système de \emph{disques} disjoints autour de ces points, dont certains
peuvent être dégénérés, réduits à leur centre. Un tel système est
\emph{admissible} s'il provient ainsi d'une surface à bord.

L'admissibilité s'exprime de façon intrinsèque. Soit $\Gamma_i$ le bord d'un
disque non dégénéré $D_i$. Le groupe $\mathbb{U}$ agit de deux façons sur
$\Gamma_i$ : comme bord du disque $D_i$, par les rotations autour du centre,
et comme composante du bord de la surface hyperbolique
$X^{\wedge} \smallsetminus \bigcup_j \mathring{D}_j$, par l'abscisse
curviligne. Le système est admissible quand ces deux structures toriques
coïncident.

Pour $X^{\wedge}$ et les points $a_i$ donnés, les systèmes admissibles
dépendent de $\nu$ paramètres, qu'on voudrait voir comme des « rayons ». Or un
disque conforme $D$ autour d'un point intérieur $a$ se plonge canoniquement
dans l'espace tangent $T_{X,a}$, comme un disque euclidien.\footnote{C'est
exact, et canonique au sens fort. Si $\varphi : (\mathring{\mathbb{D}}, 0)
\to (\mathring{D}, a)$ est une représentation conforme, l'application
$d\varphi_0 \circ \varphi^{-1}$ ne dépend pas du choix de $\varphi$, qui n'est
défini qu'à une rotation près ; elle envoie $D$ sur un disque rond centré en
$0$ et sa différentielle en $a$ est l'identité.} C'est l'\emph{ombre} $o(D)$
du disque, la notation des pages 62 et 63. Pour une métrique hermitienne
donnée sur $T_{X,a}$, c'est le disque $|z| \leqslant r$, où $r$ est ce qu'on
appelle aujourd'hui le \emph{rayon conforme} de $D$ en $a$. Les disques
possibles de $T_{X,a}$ forment un torseur $R_i$ sous $\mathbb{R}^{*}_{+}$ ;
en admettant le rayon nul, on obtient $R_i^{\wedge}$, et un système de disques
a un système de rayons $r = (r_i) \in \prod_i R_i^{\wedge}$.

Son attente est alors la suivante. Un système admissible est déterminé par
ses rayons, et un système de rayons est admissible si et seulement si
$r_i < \rho_i$ pour un système fixé $\rho = (\rho_i) \in \prod_i R_i$. Il
note lui-même que cela ne s'accorde pas avec l'action de
$(\mathbb{R}_{+})^{I}$ espérée plus haut, sauf si tous les $\rho_i$ sont
infinis. C'est la révision d'une conjecture de janvier, celle de la
page 67 : voir la section VIII, où l'on montre que les ombres sont toujours
bornées, de sorte que des $\rho_i$ finis sont inévitables.

\pagerange{18}{21}
\subsection*{Le modèle des quartiers d'orange (pages 18 à 21)}

Il avoue n'avoir pas regardé la situation de près et avoir perdu un peu le
contact depuis deux ou trois mois. Ce qui est clair pour lui, c'est que la
clef est le cas des pantalons de Thurston. Si l'on numérote les trois trous
$0, 1, \infty$, la surface bouchée $X^{\wedge}$ s'identifie canoniquement à
la sphère de Riemann $\Sigma$, et la question est de comprendre tous les
systèmes admissibles de trois disques autour de $0, 1, \infty$.

Le modèle qui lui est venu est le suivant. On munit $\Sigma$ de la métrique
sphérique pour laquelle la droite projective réelle est un grand cercle et les
points $0, 1, \infty$ sont équidistants sur cet équateur.\footnote{La lettre
dit « sa métrique euclidienne usuelle ». Pour la métrique de Fubini--Study de
la coordonnée $z$, les points $0$ et $\infty$ sont antipodaux, et non
équidistants de $1$. La métrique voulue est l'unique métrique sphérique
normalisée invariante par le groupe des six homographies qui permutent
$\{0,1,\infty\}$. Elle est unique d'après la proposition de la page 36,
puisque ce groupe n'a pas de point fixe commun sur $\Sigma$. Ses pôles sont
$e^{\pm i\pi/3}$.} Les trois points sont les centres de trois « quartiers
d'orange » $Q_0, Q_1, Q_\infty$, qui forment une décomposition cellulaire de
$\Sigma$. Le bord commun de deux quartiers est un demi-grand cercle qui passe à
égale distance des deux centres, et les trois demi-cercles se rejoignent aux
deux pôles $P^{+}, P^{-}$. Chaque $Q_i$ est un disque conforme pointé ; on
peut donc parler des disques concentriques $\lambda_i Q_i$, pour
$0 < \lambda_i < 1$, définis par l'ombre.

Le modèle affirme que les systèmes admissibles sont exactement les
$(\lambda_i Q_i)$, avec $0 \leqslant \lambda_i < 1$ si l'on admet les disques
dégénérés. Il donnerait des paramètres $\lambda_i$ adaptés à la géométrie
sphérique, liés aux rayons par $r_i = \lambda_i \rho_i$, et, pour une surface
quelconque,
\[
\overline{\mathrm{TB}} \simeq \mathrm{T} \times [0,1[^{I}.
\]
Il remarque qu'on ne peut sans doute pas espérer que la longueur $l_i$ ne
dépende que de $\lambda_i$, ce qui rend peu pertinente l'intuition du début,
fondée sur les coordonnées de Thurston. Il demande à Bers si le modèle a des
chances d'être vrai ou s'il est franchement faux.\footnote{Sous la forme où la
lettre l'énonce, disques dégénérés compris, le modèle est faux ; l'argument
qui suit est le nôtre. Prenons dégénérés les disques en $1$ et en $\infty$ :
c'est le cas d'un pantalon $P$ à deux pointes et à un bord géodésique de
longueur $L$, et le modèle prédit $D_0 = \lambda_0 Q_0 \subset Q_0$.
L'intérieur de $P$ est alors conformément
$\Sigma \smallsetminus (D_0 \cup \{1,\infty\})$, qui contient le domaine
fixe $W = (\Sigma \smallsetminus \overline{Q}_0) \smallsetminus
\{1, \infty\}$. La famille des courbes fermées de $W$ qui séparent
$\overline{Q}_0$ de $\{1,\infty\}$ a une longueur extrémale finie $c$, qui
ne dépend pas de $L$ ; elle est contenue dans la famille des courbes de $P$
librement homotopes au bord, dont la longueur extrémale est donc au plus $c$.
Mais pour la métrique hyperbolique de $P$, d'aire $2\pi$, chacune de ces
courbes a une longueur au moins $L$, puisque le bord géodésique est la plus
courte de sa classe : la longueur extrémale est au moins $L^{2}/2\pi$. Pour
$L^{2} > 2\pi c$, c'est contradictoire. Le disque admissible $D_0$ déborde
donc de $Q_0$ dès que le bord est assez long. Le dossier ne contient pas de
réponse de Bers.}

La lettre se termine par une difficulté plus générale. En dehors de
l'existence et de l'unicité de la métrique hyperbolique, compatible à la
structure conforme et pour laquelle le bord est géodésique, on a peu de prise
sur ses propriétés. Que devient par exemple la métrique d'un pantalon quand
on lui retire un collier ouvert le long du bord ? Il souhaiterait, si le
modèle est juste, une expression explicite de la métrique d'un pantalon en
fonction des $\lambda_i$.

\pagerange{23}{45}
\section*{IV. Métriques canoniques sur les surfaces conformes (pages 23 à 45)}

La chemise de la page 23 porte le titre « Métriques canoniques sur les
surfaces conformes ». Les feuillets se lisent dans l'ordre établi plus haut :
26, 44, 27, 45, 28 à 43. Deux feuillets volants les précèdent.

\pagerange{24}{25}
\subsection*{Données équivalentes sur une sphère conforme (pages 24 et 25)}

Sur une sphère conforme $X$, de groupe $G = \mathrm{Aut}^{+}(X)$, la page 24
donne des données équivalentes :

\begin{itemize}
  \item[a)] une paire compatible formée d'un antipodisme et d'une
  inversion, c'est-à-dire deux structures réelles compatibles ;
  \item[b)] une métrique sphérique et un grand cercle ;
  \item[c)] un cercle muni de l'inversion qui le fixe, et une paire de points
  échangés par elle ;
  \item[d)] un tore compact $T \subset G$ et une orbite de $T$ dans
  $X \smallsetminus X^{T}$ ;
  \item[d')] une paire de points $\{z, z'\}$ et une orbite, dans
  $X \smallsetminus \{z,z'\}$, de l'unique tore compact qui les fixe.
\end{itemize}

Tout se vérifie en plaçant les pôles en $0$ et $\infty$. La métrique
sphérique d'équateur $|z| = r$ et de pôles $0, \infty$ est unique, et le
cercle $|z| = r$ est une orbite du tore $\{z \mapsto tz, |t| = 1\}$. La
page 25 est une feuille de calculs isolés ; on y lit déjà la suite exacte
$0 \to \Pi \to V \to E \to 0$ et l'isomorphisme
$H^{1}(\gamma, E) \simeq H^{2}(\gamma, \Pi) \simeq \Pi^{\gamma}/(1 +
\sigma)\Pi$ qui clôt la section XII.

\pagerange{26}{44}
\subsection*{Surfaces hyperboliques (pages 26 et 44)}

Soit $X$ une surface holomorphe connexe (séparée), non nécessairement
compacte, et $\widetilde{X}$ son revêtement universel.

\emph{Théorème.} Les conditions suivantes sont équivalentes :
\begin{itemize}
  \item[(i)] $X$ n'est isomorphe ni à $\mathbb{P}^{1}_{\mathbb{C}}$, ni à
  $\mathbb{C}$, ni à $\mathbb{C}^{*}$, ni à une courbe elliptique ;
  \item[(ii)] $\widetilde{X}$ est isomorphe au demi-plan $\mathbb{H}$ ;
  \item[(iii)] toute application holomorphe $\mathbb{C} \to X$ est constante ;
  \item[(iv)] si $X$ est un ouvert d'une droite projective $X'$ (resp.\ d'une
  droite affine $X'$), le complémentaire $X' \smallsetminus X$ a au moins
  trois (resp.\ deux) points ;
  \item[(v)] si $X = \widehat{X} \smallsetminus S$, avec $\widehat{X}$
  compacte de genre $g$ et $S$ finie de cardinal $\nu$, on a
  $2g + \nu \geqslant 3$, soit $(g,\nu) \notin \{(0,0),(0,1),(0,2),(1,0)\}$,
  soit $\chi(X) < 0$.
\end{itemize}
On dit alors que $X$ est \emph{hyperbolique}.

C'est le théorème d'uniformisation de Koebe et Poincaré, avec le théorème de
Liouville pour (iii). Deux corollaires suivent. \emph{Corollaire 1} : un ouvert
connexe d'une surface hyperbolique est hyperbolique, par (iii).
\emph{Corollaire 2} : si $X$ est une surface holomorphe connexe à bord non
vide, $\operatorname{Int} X$ est hyperbolique.\footnote{La page ne le démontre
pas. On le voit sur le double $X^{\mathrm{or}}$ : $\operatorname{Int} X$ y
est un ouvert dont le complémentaire contient un arc. Si $X^{\mathrm{or}}$ est
hyperbolique, on conclut par le corollaire 1 ; si c'est $\mathbb{P}^{1}$,
$\mathbb{C}$ ou $\mathbb{C}^{*}$, par (iv) ; si c'est une courbe elliptique,
parce que $\operatorname{Int} X$ est contenu dans la courbe privée d'un
point.}

La page 44 ajoute un \emph{Corollaire 3}. Soit $X'$ une surface holomorphe
connexe et $X \subset X'$ un ouvert connexe. Alors $X$ est hyperbolique, sauf
si $X'$ est isomorphe à $\mathbb{P}^{1}$ (resp.\ $\mathbb{C}$, resp.\
$\mathbb{C}^{*}$) et $X' \smallsetminus X$ a au plus deux points (resp.\ au
plus un, resp.\ aucun), ou si $X'$ est compacte de genre $1$ et
$X = X'$.\footnote{La page écrit $\operatorname{card}(X \smallsetminus X')$ ;
il faut $X' \smallsetminus X$.}

Elle ajoute aussi une \emph{Remarque} : une surface conforme connexe $X$,
orientable ou non, est dite hyperbolique si $\widetilde{X} \simeq \mathbb{H}$,
ou, ce qui revient au même, si toute application conforme $\mathbb{C} \to X$
est constante, ou encore si son revêtement des orientations est hyperbolique.
Les surfaces conformes connexes sans bord non hyperboliques sont donc les
quatre surfaces de (i), plus leurs quotients par des anti-involutions sans
point fixe.

\pagerange{27}{30}
\subsection*{Les surfaces non hyperboliques, et leurs métriques (pages 27, 45, 28 à 30)}

Ces quotients sont les suivants (page 27).

\begin{itemize}
  \item[a)] Le plan projectif réel, quotient de la sphère par l'antipodie, et
  le plan projectif réel privé d'un point, quotient de $\mathbb{C}^{*}$ par
  $z \mapsto -1/\bar z$. Leurs revêtements des orientations sont
  respectivement $\mathbb{P}^{1}$ et $\mathbb{C}^{*}$.\footnote{La page écrit
  « revêtements universels ». Le revêtement universel du plan projectif réel
  privé d'un point, qui est un ruban de Möbius ouvert, est $\mathbb{C}$ et non
  $\mathbb{C}^{*}$ ; c'est bien le revêtement des orientations qui est en
  jeu, comme le dit la page 28 en écrivant $X^{\mathrm{or}} \simeq
  \mathbb{C}^{*}$.} Aucune anti-involution de $\mathbb{C}$ n'est sans point
  fixe (page 42), de sorte que $\mathbb{C}$ ne donne rien.
  \item[b)] Les surfaces $E(\mathbb{C})/\text{conj}$, où $E$ est une courbe
  réelle lisse géométriquement connexe de genre $1$ sans point réel : ce sont
  des bouteilles de Klein. La page les nomme $E_a$, $a \in \mathbb{R}^{*}_{+}$ ;
  la page 100 les identifie aux torseurs non triviaux $E_a^{!}$ sous les
  courbes $E_a$ de la section XII.
\end{itemize}

Suit la question des métriques riemanniennes canoniques, compatibles à la
structure conforme, sur une surface conforme sans bord.

\begin{enumerate}
  \item \emph{Cas hyperbolique.} La métrique de Poincaré de $\widetilde{X}
  \simeq \mathbb{H}$, de courbure $-1$, est invariante par tous les
  automorphismes conformes de $\mathbb{H}$, antiholomorphes compris ; elle
  descend à $X$. Si $X$ est compacte, c'est l'unique métrique compatible de
  courbure constante $-1$ (page 45).
  \item \emph{Genre $1$, compact} (pages 45 et 28). Sur une courbe elliptique,
  les métriques plates invariantes par translation forment une classe à un
  scalaire près. On la normalise en imposant l'aire totale $1$. La métrique
  obtenue est invariante par tous les automorphismes conformes, et elle
  descend donc aux bouteilles de Klein.
  \item[2')] \emph{$X \simeq \mathbb{C}^{*}$, ou plan projectif réel privé d'un
  point.} Les automorphismes holomorphes de $\mathbb{C}^{*}$ sont
  $z \mapsto az$ et $z \mapsto a/z$. Ils relèvent des translations et des
  symétries de $\mathbb{C}$ par $z \mapsto \exp 2i\pi z$, et conservent la
  métrique plate $|dz|^{2}/|z|^{2}$. On la normalise en donnant au cercle
  $|z| = 1$ la longueur $1$, ce qui fait
  $|dz|^{2}/4\pi^{2}|z|^{2}$, l'image de la métrique $|dz|^{2}$ de
  $\mathbb{C}$ par $z \mapsto \exp 2i\pi z$, comme le dit la marge. La
  conjugaison la conserve aussi, et la métrique descend au plan projectif
  privé d'un point.
  \item[2'')] \emph{$X \simeq \mathbb{C}$.} Les métriques invariantes par
  translation correspondent aux formes hermitiennes définies positives sur
  l'espace des translations, définies à homothétie près. Elles sont
  invariantes par les translations et les rotations, et seulement
  quasi-invariantes, c'est-à-dire à un facteur constant près, par les
  automorphismes affines quelconques. On n'a donc qu'une classe de métriques
  à un scalaire positif près.\footnote{La page intitule ce cas « $X \simeq
  \mathbb{E}_{\mathbb{C}}$ » et parle de métriques sur une courbe
  « elliptique » (lecture incertaine) complexe. Tout ce qui est dit ---
  automorphismes affines, structure canonique de droite affine sur
  $\mathbb{C}$, classe définie seulement à un scalaire près --- concerne la
  droite affine $\mathbb{E}^{1}_{\mathbb{C}} \simeq \mathbb{C}$, et c'est le
  cas a) de la récapitulation de la page 37. Sur une courbe elliptique, au
  contraire, la normalisation par l'aire lève le scalaire (cas 2).}
  \item[3)] \emph{Plan projectif réel.} $X^{\mathrm{or}}$ est une droite
  projective munie d'un antipodisme, ce qui équivaut à une métrique sphérique
  compatible, normalisée par l'aire $1$. Cette métrique descend à $X$.
  \item[3')] \emph{$X \simeq \mathbb{P}^{1}_{\mathbb{C}}$.} Aucune métrique
  n'est invariante, ni même quasi-invariante, par tous les automorphismes. Il
  n'y a donc pas de métrique canonique sans structure supplémentaire.
\end{enumerate}

\pagerange{30}{36}
\subsection*{La sphère munie d'un groupe compact (pages 30 à 36)}

Soit $X = \mathbb{P}(V)$, où $V$ est un plan vectoriel complexe, et
$G = \mathrm{Aut}_{\mathrm{conf}}(X)$, antiholomorphes compris. On se donne
un sous-groupe $H \subset G$. Pour qu'il laisse invariante une métrique
sphérique, il faut, s'il est fermé, qu'il soit compact. Soit
$H^{+} = H \cap \mathrm{Aut}^{+}(X)$. Son image réciproque dans
$\mathrm{SL}(V)$ est compacte ; elle laisse invariante une forme hermitienne
définie positive sur $V$, d'où une métrique sphérique invariante, normalisée
par $\int_X dx = 1$. Si la forme invariante est unique à un scalaire près, la
métrique est canonique : elle ne dépend que de $H^{+}$ et elle est invariante
par le normalisateur de $H^{+}$, a fortiori par $H$.

La forme invariante n'est pas unique exactement quand $H^{+}$ fixe deux points
distincts $z, z'$. Alors $H^{+}$ est contenu dans le stabilisateur de $z$ et
$z'$, isomorphe à $\mathbb{C}^{*}$, et dans son unique sous-groupe compact
maximal $T \simeq \mathbb{U}$ ; si $H^{+} \neq \{1\}$, la paire $\{z,z'\}$
est déterminée par $H^{+}$.\footnote{La marge dit « à vérifier ». C'est le
lemme de Schur : l'image réciproque de $H^{+}$ agit sur $V$ de façon
complètement réductible, et la forme invariante est unique à un scalaire près
si et seulement si $V$ est irréductible.} Si $H^{+} \neq \{1\}$, les
métriques invariantes par $H^{+}$ sont celles qui sont invariantes par $T$.
Elles correspondent bijectivement aux orbites de $T$ dans $X \smallsetminus
\{z, z'\}$, sur lequel $T$ agit librement : on associe à une métrique son
équateur, le grand cercle équidistant de $z$ et $z'$.

Si $H \neq H^{+}$, deux cas se présentent.

\emph{a)} $H^{+} = \{1\}$ et $H = \{1, \sigma\}$, avec $\sigma$ une
anti-involution. Si $\sigma$ est sans point fixe (un antipodisme), elle fixe
une unique métrique sphérique normalisée. Si $\sigma$ a des points fixes (une
inversion), les métriques sphériques invariantes correspondent aux paires
$\{x, \sigma x\}$ de points échangés, qui en sont les pôles, c'est-à-dire aux
points de $(X \smallsetminus X^{\sigma})/\sigma$.

\emph{b)} $H^{+} \neq \{1\}$. Soit $T$ le tore compact maximal qui contient
$H^{+}$, et $N = T\cdot H$, extension de $\gamma = \mathbb{Z}/2\mathbb{Z}$ par
$T$. Un élément antiholomorphe de $N$ qui \emph{fixe} $z$ et $z'$ agit sur $T$
par $t \mapsto t^{-1}$ ; un élément qui les \emph{échange} centralise
$T$.\footnote{La page écrit d'abord l'inverse, puis porte en marge « c'est
l'inverse ! ». Le calcul est immédiat : en plaçant $\{z,z'\}$ en
$\{0,\infty\}$, si $\sigma(z) = a\bar z$ et $t(z) = \theta z$, on trouve
$\sigma t \sigma^{-1}(z) = \bar\theta z$ ; si $\sigma(z) = b/\bar z$, on
trouve $\sigma t\sigma^{-1} = t$. C'est la version de la marge qu'on suit,
et la numérotation des cas qui suit est rectifiée en conséquence.}

\begin{itemize}
  \item[1)] \emph{$N$ extension centrale.} Comme $H^{2}(\gamma, T) =
  T/T^{2} = 0$ pour l'action triviale, $N \simeq T \times \{1, \sigma\}$, où
  $\sigma$ échange $z$ et $z'$. Avec $\{z,z'\} = \{0, \infty\}$, les
  anti-automorphismes qui échangent $0$ et $\infty$ sont les
  $z \mapsto b/\bar z$, et ce sont eux qui centralisent $T$.\footnote{La page
  dit « normalisent » ; les deux familles normalisent $T$, et c'est la
  seconde qui le centralise.} On a $\sigma^{2}(z) = (b/\bar b)\,z$ : $\sigma$
  est involutive si et seulement si $b$ est réel. Ses points fixes sont donnés
  par $z\bar z = b$ : c'est une inversion si $b > 0$ et un antipodisme si
  $b < 0$. On peut remplacer $\sigma$ par $\sigma' = \varepsilon\sigma$, où
  $\varepsilon = -1$ est l'élément d'ordre $2$ de $T$ : des deux, l'un est
  une inversion et l'autre un antipodisme. Si $b > 0$, l'inversion $\sigma$
  fixe le cercle $C = \{z\bar z = b\}$, et l'antipodisme
  $\sigma'(z) = -b/\bar z$ conserve $C$ : tous deux définissent la même
  métrique, celle de pôles $0, \infty$ et d'équateur $C$. « On est
  content ! »
  \item[2)] \emph{$N$ non centrale.} Alors $\sigma(z) = a\bar z$, et
  $\sigma^{2} = (z \mapsto a\bar a z)$ appartient à $T$ si et seulement si
  $a \in \mathbb{U}$, auquel cas $\sigma^{2} = 1$ et l'extension est scindée.
  Si $a = \alpha^{2}$, avec $\alpha \in \mathbb{U}$, les points fixes de
  $\sigma$ forment la droite $\alpha\mathbb{R}$, un méridien. La marge
  conclut : $N = T \cup N^{-}$, où $N^{-}$ est formé de toutes les
  anti-involutions qui fixent $z$ et $z'$, et $N$ ne détermine pas plus de
  métrique que $T$.
\end{itemize}

\emph{Proposition} (page 36, « en résumé et sauf erreur »). Soit $X$ une
sphère conforme et $H \subset G$ un sous-groupe compact.
\begin{itemize}
  \item[a)] Il existe sur $X$ une métrique sphérique normalisée, compatible,
  invariante par $H$.
  \item[b)] Elle est unique si et seulement si $X^{H} = \emptyset$.
  \item[c)] Si $X^{H} \neq \emptyset$ et $H \neq \{1\}$, alors ou bien
  $\operatorname{card} X^{H} = 2$, ou bien $H = \{1,\sigma\}$ avec $\sigma$
  une inversion. Dans le premier cas, $X^{H} = X^{H^{+}}$, $H^{+}$ est
  contenu dans un unique tore compact $T$, et les métriques invariantes par
  $H$ correspondent bijectivement aux orbites de $T$ dans
  $X \smallsetminus X^{H}$. Dans le second, elles correspondent aux éléments de
  $(X \smallsetminus X^{\sigma})/\sigma$.
\end{itemize}

La marge demande de traduire l'action de $G$ sur les métriques sphériques en
termes d'algèbre linéaire. La traduction est la suivante, et elle est la
nôtre. Les métriques sphériques normalisées sont les formes hermitiennes
définies positives sur $V$ à un scalaire près, et elles forment l'espace
hyperbolique $\mathbb{H}^{3} = \mathrm{PSL}_2(\mathbb{C})/\mathrm{PSU}(2)$,
sur lequel $G$ agit par isométries, les antiholomorphes renversant
l'orientation. Les métriques invariantes par $H$ forment l'ensemble des
points fixes de $H$ dans $\mathbb{H}^{3}$, qui est une sous-variété totalement
géodésique non vide. La proposition dit laquelle : un point si
$X^{H} = \emptyset$ ; la géodésique d'extrémités $z, z'$ si
$\operatorname{card} X^{H} = 2$, paramétrée par les orbites de $T$ ; le plan
hyperbolique bordé par le cercle $X^{\sigma}$ si $\sigma$ est une inversion.
C'est une surface isomorphe au disque $(X \smallsetminus X^{\sigma})/\sigma$.

\pagerange{37}{38}
\subsection*{Récapitulation (pages 37 et 38)}

Sur toute surface conforme sans bord $X$, on peut définir une métrique
riemannienne canonique $Q_X$, compatible à la structure conforme, sauf dans
deux cas.

\begin{itemize}
  \item[a)] $X \simeq \mathbb{C}$. On n'a qu'une classe de métriques à un
  scalaire près : les métriques invariantes par translation, quasi-invariantes
  par tout automorphisme conforme, et plates.
  \item[b)] $X \simeq \mathbb{P}^{1}_{\mathbb{C}}$. Aucune métrique n'est
  quasi-invariante par $G$, ni même par $\mathrm{Aut}^{+}(X)$. Les métriques
  sphériques normalisées forment une orbite $R(X)$ de $G$, variété de
  dimension $3$. Il existe un $Q \in R(X)$ invariant par $H \subset G$ si et
  seulement si $H$ est relativement compact, et il est unique si et seulement
  si $X^{H} = \emptyset$.
\end{itemize}

Quand $Q_X$ est défini, sa courbure est constante. Elle vaut $-1$ si et
seulement si $X$ est hyperbolique ; elle est nulle si et seulement si
$\widetilde{X} \simeq \mathbb{C}$, cas a) compris ; elle est positive dans le
seul cas restant, celui du plan projectif réel. Il en va de même de la
métrique $Q_{X,H}$ définie par un groupe compact $H$ sans point fixe sur
$\mathbb{P}^{1}$ : elle est de courbure constante positive. La marge
généralise : pour toute surface conforme connexe $X$, l'ensemble des métriques
compatibles contient une orbite canonique de $\mathrm{Aut}_{\mathrm{conf}}(X)$,
réduite à un point sauf dans les cas a) et b), où c'est une demi-droite et
une variété de dimension $3$. Elle propose le nom de \emph{métriques
distinguées}.\footnote{Un alinéa entre crochets, en tête de la page 38,
envisage le cas où $X$ est associé à une conique et $H$ au groupe de ses
automorphismes ; il est trop lacunaire pour être restitué.}

\pagerange{39}{43}
\subsection*{Surfaces à bord (pages 39 à 43)}

Le cas des surfaces à bord se ramène au précédent par l'équivalence de
catégories
\[
X \longmapsto (X^{\mathrm{or}}, \sigma), \qquad
(X', \sigma) \longmapsto X'/\sigma,
\]
entre surfaces conformes à bord et surfaces holomorphes sans bord munies d'une
anti-involution. C'est aujourd'hui la théorie des surfaces de Klein, et
$X^{\mathrm{or}}$ est le double complexe : le revêtement des orientations,
replié le long du bord.\footnote{La page note ce double par un exposant lu
« or », et la page 41 par $\mathfrak{X}$. L'équivalence vaut pour les
morphismes qui envoient le bord dans le bord, que la page ne précise pas. La
théorie est exposée dans Alling et Greenleaf, \emph{Foundations of the theory
of Klein surfaces} (1971), que la page ne cite pas.} On dit que $X$ est
hyperbolique si $X^{\mathrm{or}}$ l'est. Pour $X$ sans bord, cela redonne la
définition précédente, car $X^{\mathrm{or}}$ est alors un revêtement étale de
$X$. Pour $X$ compacte, on a $\chi(X^{\mathrm{or}}) = 2\chi(X)$, puisque le
bord, réunion de cercles, a une caractéristique nulle ; donc $X$ est
hyperbolique si et seulement si $\chi(X) < 0$. Si $X$ est orientable, de genre
$g$ à $\nu$ bords, c'est $2g + \nu \geqslant 3$. Il faut se garder de
confondre avec l'hyperbolicité de $\operatorname{Int} X$, qui a toujours lieu
sauf si $\nu = 0$ et $g \in \{0,1\}$ : le disque fermé, par exemple, n'est pas
hyperbolique, alors que son intérieur l'est.

Si $X$ est hyperbolique, la métrique hyperbolique de $X^{\mathrm{or}}$ est
invariante par $\sigma$, et elle descend en une métrique hyperbolique sur $X$
pour laquelle le bord, lieu des points fixes d'une isométrie involutive, est
géodésique. Plus généralement, toute métrique canonique de $X^{\mathrm{or}}$
invariante par $\sigma$ descend, avec un bord géodésique. Pour le plan
projectif réel, il faut exiger que la métrique de $X^{\mathrm{or}}$ soit dans
$R(X^{\mathrm{or}})$ et stable par $\sigma$, ce qui la détermine. Pour $X
\simeq \mathbb{C}$, on a $X^{\mathrm{or}} \simeq \mathbb{C} \amalg
\mathbb{C}$, et la prescription ne fournit une métrique que si on la complète
par l'hypothèse qu'on puisse choisir sur $X^{\mathrm{or}}$ une métrique
distinguée stable par $\sigma$ (bas de la page 40 et haut de la page 41).

Restent les surfaces connexes, compactes ou non, à bord non vide : un double
connexe $X^{\mathrm{or}}$ et une anti-involution $\sigma$ à points fixes. Les
seuls cas exceptionnels sont les suivants.

\begin{itemize}
  \item[a')] $X^{\mathrm{or}} \simeq \mathbb{C}$. Les anti-automorphismes de
  $\mathbb{C}$ sont les $z \mapsto a\bar z + b$. Une anti-involution a
  toujours un point fixe ; si on le place en $0$, elle s'écrit
  $\sigma(z) = \alpha^{2}\bar z$ avec $\alpha \in \mathbb{U}$, et c'est la
  symétrie orthogonale par rapport à la droite $\alpha\mathbb{R}$. La surface
  $X$ est un demi-plan fermé $\overline{\mathbb{H}}$. Son groupe
  d'automorphismes conformes, celui des automorphismes de $\mathbb{C}$ qui
  commutent à la conjugaison modulo celle-ci, est le groupe affine réel. Les
  métriques plates invariantes par translation, invariantes par les
  automorphismes d'ordre fini et en particulier par $\sigma$, donnent sur
  $X$ une métrique canonique à une constante près.
  \item[b')] $X^{\mathrm{or}} \simeq \mathbb{P}^{1}$, $\sigma$ une inversion,
  qui définit une structure de droite projective réelle. $X$ est un disque
  fermé, et une métrique distinguée invariante par $\sigma$ demande le choix
  d'un point du disque ouvert.
\end{itemize}

\emph{En résumé} (page 43), une surface conforme à bord connexe a une unique
métrique distinguée, sauf dans quatre cas : a) $X \simeq \mathbb{C}$ et a')
$X \simeq \overline{\mathbb{H}}$, où l'on a une classe de métriques à une
constante près, orbite sous $\mathrm{Aut}\,X$ ; b)
$X \simeq \mathbb{P}^{1}$ et b') $X$ un disque, où les métriques distinguées
forment une orbite de dimension $3$, respectivement $2$.

\pagerange{47}{50}
\section*{V. Isomorphismes de Teichmüller et questions anabéliennes (pages 47 à 50)}

Le feuillet de la page 47 est daté « 15.1. », sans année. On se donne une
courbe $X$ lisse de type $(g,\nu)$ sur $\overline{\mathbb{Q}}$, et une courbe
$Y$ stable non lisse de même type, munie d'un « propagateur normal » $N$ en
ses points doubles. On obtient un groupe fondamental $\pi_1([Y,N])$ et une
notion d'\emph{isomorphisme de Teichmüller} $X \simeq_T Y$. Un tel isomorphisme
induit un isomorphisme de groupes extérieurs $\pi_1(X) \to \pi_1([Y,N])$, et
il est déterminé par lui. Avec des points base, ordinaires ou à l'infini, on
obtient un isomorphisme de groupes. Cet isomorphisme n'est pas compatible aux
actions extérieures de $\mathrm{Gal}(\overline{\mathbb{Q}}/K)$, même pour $K$
assez grand.\footnote{Le propagateur normal est défini à la page 74 : en
chaque point double, dont les deux branches passent par $x_1$ et $x_2$ sur la
normalisée, un élément non nul de $T_{x_1} \otimes T_{x_2}$. C'est un vecteur
normal non nul à la strate des courbes à point double, puisque le paramètre de
lissage d'un nœud vit dans cette droite. $\pi_1([Y,N])$ se lit donc comme le
groupe fondamental pris en un point base tangentiel. Cette notion a été
dégagée plus tard par Deligne (1989), et utilisée pour les groupes de
Teichmüller par Ihara et Nakamura (1997) ; ce sont des repères postérieurs,
que la page ne cite pas.} Il note qu'il faudrait mettre en évidence trois
classes de conjugaison d'homomorphismes $\hat{\mathbb{Z}}(1) \to \pi_1(X)$ ---
deux seulement si $\nu = 0$, l'une des trois correspondant aux bouts ---,
sans dire lesquelles.

\emph{À examiner} (page 48) :
\begin{itemize}
  \item[a)] étendre la « conjecture fondamentale » de la géométrie anabélienne
  aux courbes anabéliennes singulières, stables et munies d'un propagateur
  normal ;
  \item[b)] pour une telle courbe $Y$ sur un corps $K$ de type fini sur
  $\mathbb{Q}$, récupérer les points singuliers et les points marqués à partir
  des classes de conjugaison de sous-groupes de $\pi_1(Y)$ isomorphes à
  $\hat{\mathbb{Z}}$ et normalisés par $\pi_1(K)$, qui agit alors sur eux par
  le caractère cyclotomique ;
  \item[c)] à l'image des compactifications de Deligne et Mumford des
  $M_{g,\nu}$, définir une compactification anabélienne d'une variété lisse
  anabélienne, et étudier son comportement par spécialisation.
\end{itemize}

La conjecture fondamentale est celle que Grothendieck avait formulée dans sa
lettre à Faltings du 27 juin 1983, pour les courbes hyperboliques sur les
corps de type fini sur $\mathbb{Q}$.\footnote{« MD », dans c), est l'abréviation
de Mumford--Deligne, qu'il écrit en toutes lettres à la page 75. La conjecture
pour les courbes hyperboliques a été démontrée par Tamagawa (cas affine, 1997)
et Mochizuki (1999). La reconnaissance des sous-groupes d'inertie des pointes,
qui est la question b) pour une courbe lisse, remonte à Nakamura (1990) ; la
géométrie anabélienne des courbes stables a été développée ensuite, notamment
par Mochizuki. Tout cela est postérieur à la page et n'y est évidemment pas
cité.}

La page 49 attend trois types de présentation par générateurs et relations
pour les groupoïdes et groupes de Teichmüller, chacun à signification
intrinsèque.

\begin{itemize}
  \item[a)] Par \emph{pantalons}, ou triades. Elle revient à travailler avec
  les points à l'infini les plus dégénérés des multiplicités modulaires, et
  elle semble commode pour les groupes purs $\Gamma^{+!}_{g,\nu}$, d'après
  l'exemple de $\Gamma^{+!}_{0,4}$.\footnote{La page dit « les plus
  antiques ». Les points les plus dégénérés sont les courbes stables dont
  toutes les composantes sont des droites projectives à trois points spéciaux.
  Hatcher et Thurston avaient donné en 1980 une présentation du groupe
  modulaire fondée sur les décompositions en pantalons.}
  \item[b)] Par une « boîte » complète de pièces de Lego : les groupes
  d'automorphismes des pièces, et certains chemins de déformation entre
  pièces. Exemples frappants : $\Gamma^{+}_{0,4}$, et
  $\Gamma^{+}_{1,1} \simeq \mathrm{SL}_2(\mathbb{Z})$, engendré par $\rho$ et
  $\sigma$ avec $\rho^{3} = \sigma^{2}$ et $\rho^{6} = \sigma^{4} = 1$. C'est
  la présentation classique $\mathrm{SL}_2(\mathbb{Z}) \simeq \mathbb{Z}/6
  *_{\mathbb{Z}/2} \mathbb{Z}/4$.
  \item[c)] Par \emph{involutions}, qui sont dans $\Gamma^{-}_{g,\nu}$.
  Exemple typique : $\Gamma_{1,1} \simeq \mathrm{GL}_2(\mathbb{Z})$, engendré
  par trois réflexions $r_1, r_2, r_3$ avec
  \[
  (r_1r_2)^{2} = (r_2r_3)^{3}, \qquad (r_1r_2)^{4} = 1 .
  \]
  Avec les relations $r_i^{2} = 1$, implicites dans le mot « réflexions »,
  c'est la présentation de $\mathrm{GL}_2(\mathbb{Z})$ comme produit amalgamé
  $D_4 *_{D_2} D_6$.\footnote{Les exposants de la dernière relation sont
  surchargés sur la page ; la lecture donnée est la dernière encre, et c'est
  celle qui fait de $\langle r_1,r_2\rangle$ un groupe diédral d'ordre $8$ et
  de $\langle r_2,r_3\rangle$ un groupe diédral d'ordre $12$, amalgamés le
  long de $\langle r_2, (r_1r_2)^{2}\rangle$.}
\end{itemize}

Le « jeu de Lego » et l'idée que les groupes de Teichmüller se présentent à
partir des petites dimensions sont ceux de l'\emph{Esquisse d'un programme},
rédigée ce même mois de janvier 1984.\footnote{La question de la tour de
Teichmüller a été reprise par Hatcher, Lochak et Schneps (2000). C'est un
contexte postérieur.}

La page 50 est un brouillon écrit en tous sens. On y lit la suite exacte
\[
1 \to \Pi_{0,3} \to \Gamma^{+}_{0,4} \to \mathfrak{S}_4 \to 1,
\]
où $\Pi_{0,3}$, groupe fondamental de la sphère privée de trois points, est
libre de rang $2$ et s'identifie au groupe pur $\Gamma^{+!}_{0,4}$. Elle
s'envoie sur la suite
$1 \to \Pi_{0,3} \to \Gamma^{+}_{1,1}/\pm 1 \to \mathfrak{S}_3 \to 1$, avec
$\Gamma^{+}_{1,1}/\pm 1 \simeq \mathrm{PSL}_2(\mathbb{Z})$ et $\mathfrak{S}_3
\simeq \mathrm{SL}_2(\mathbb{Z}/2)$.\footnote{La page écrit l'indice
« $0,1$ » ; l'isomorphisme avec $\mathrm{SL}_2(\mathbb{Z})/\pm 1$ demande
$(1,1)$, comme la transcription le signale. Le noyau de
$\mathfrak{S}_4 \to \mathfrak{S}_3$ est le groupe de Klein, qui est l'objet
de la section XI.} On y lit aussi les relations entre lacets autour des
pointes, comme $\ell_0\ell_1\ell_\infty = 1$ ; l'identification
$\mathrm{SL}_2(\mathbb{Z})/\pm B \simeq \mathbb{P}^{1}(\mathbb{Q})$, où $B$
est le groupe unipotent, les cusps de la courbe modulaire ; la suite
$1 \to \Pi_{g,\nu-1} \to \Gamma_{g,\nu} \to \Gamma_{g,\nu-1} \to 1$ de l'oubli
d'un point ; et une liste de notations (groupoïdes, groupes, multiplicités,
espace de Teichmüller).

\pagerange{52}{52}
\section*{VI. Un brouillon : deux disques et une dualité (page 52)}

Page très surchargée. On y lit d'abord des correspondances. Un cercle torique
orienté est un torseur sous $\mathbb{U}$. Les droites vectorielles
complexes unitaires, c'est-à-dire munies d'une norme hermitienne, en donnent
par leur cercle unité ; pour une telle droite $T$, on a $T^{\vee} \simeq
\overline{T}$. Puis vient la correspondance qui servira à la page 74. Soient
$T'$ et $T''$ deux droites vectorielles complexes. La donnée de disques
$D' \subset T'$ et $D'' \subset T''$ et d'un anti-isomorphisme $u$ de
$\mathbb{U}$-torseurs entre leurs bords équivaut à la donnée d'une dualité
$v : T' \otimes T'' \to \mathbb{C}$ et du seul disque $D'$. On retrouve
$D''$ comme disque polaire de $D'$ pour $v$, et $u$ comme l'identification
des bords qu'impose $v$. Le compte est bon : deux rayons et un angle d'un
côté, un rayon et un nombre complexe non nul de l'autre.\footnote{Un mot
abrégé, lu « Bi… », suggère les formes bilinéaires ; l'identification
de la fin de la chaîne à $T'^{\times} \wedge_{\mathbb{C}^{\times}}
T''^{\times}$ est celle de la page 74. La vérification du compte est la
nôtre.}

\pagerange{53}{60}
\section*{VII. Chirurgie conforme : les fondements (pages 53 à 60)}

Le feuillet de la page 53 porte « 11.1.84 », « 1 » et le titre souligné
« Chirurgie conforme ». La suite, qui va jusqu'à la page 81, est numérotée
par lui en tête des rectos : 1 à 4 aux pages 53, 55, 57 et 59, puis, le
chiffre de la page 60 étant raturé, 5 à 15 aux pages 61 à 81.

\pagerange{53}{55}
\subsection*{(1) Surfaces conformes à bord (pages 53 à 55)}

Une surface conforme à bord se définit par des cartes modelées, au voisinage
du bord, sur $\overline{\mathbb{H}}$.\footnote{Le numéro cerclé en tête peut
se lire « 4 » ; la suite (2), (3), (4) des pages 55, 57 et 60 demande (1).}
Dans le cas orienté, la structure est un faisceau de $\mathbb{C}$-algèbres :
en un point $x \in \mathbb{R} = \partial\overline{\mathbb{H}}$, ce sont les
germes de fonctions qui se \emph{prolongent} en fonctions holomorphes au
voisinage de $x$ dans $\mathbb{C}$. Ce faisceau est plus petit que celui des
fonctions continues holomorphes à l'intérieur, et il ne faut pas les
confondre.

\emph{Théorème 1.} Soit $X$ une surface topologique à bord. Une structure
conforme à bord sur $X$ est déterminée par sa restriction à
$\operatorname{Int} X$.\footnote{La page dit que la structure de $X$ « est
connue quand on connaît celle de $\operatorname{Int}(X)$ » ; il faut entendre
avec la topologie de $X$, qu'aucune structure de l'intérieur ne détermine
(le disque ouvert est l'intérieur du disque fermé comme de bien d'autres
surfaces à bord). C'est un énoncé d'unicité, non d'existence.} Les cartes au
bord sont exactement les plongements topologiques ouverts
$U \hookrightarrow \overline{\mathbb{H}}$, qui envoient bord sur bord et sont
holomorphes à l'intérieur. Cela résulte du

\emph{Lemme 1.} Soit $U$ un ouvert de $\overline{\mathbb{H}}$ et
$f : U \hookrightarrow \overline{\mathbb{H}}$ un plongement ouvert,
holomorphe dans $U \cap \mathbb{H}$. Alors $f$ se prolonge holomorphiquement
au voisinage de $U \cap \mathbb{R}$ dans $\mathbb{C}$, et sa dérivée ne s'y
annule pas.

La première assertion est le principe de réflexion de Schwarz. Une fonction
continue sur $U$, holomorphe dans $U \cap \mathbb{H}$ et réelle sur
$U \cap \mathbb{R}$, se prolonge par $f(\bar z) = \overline{f(z)}$ ; et $f$
est réelle sur le bord, puisqu'un plongement ouvert envoie les points du bord
sur le bord.\footnote{La page dit « lemme de Schwarz » et omet l'hypothèse
d'holomorphie à l'intérieur, sans laquelle l'énoncé est faux. Le principe de
réflexion date de 1870.} La seconde vient de ce que, si le développement en
$x = 0$ commence par $a_n z^{n}$ avec $a_n \neq 0$, l'injectivité et
l'inclusion dans $\overline{\mathbb{H}}$ imposent $n = 1$.

Un énoncé est mis entre crochets : le \emph{théorème 2}, équivalence entre
surfaces conformes compactes à bord et courbes algébriques propres et lisses
sur $\mathbb{R}$. C'est la version algébrique de l'équivalence de la
page 39 : $X \mapsto$ la courbe réelle dont $X^{\mathrm{or}}$ est la
surface de Riemann et $\sigma$ la conjugaison.\footnote{Pour les morphismes
qui envoient le bord dans le bord. Les courbes réelles non géométriquement
connexes correspondent aux surfaces orientables sans bord. C'est la théorie
d'Alling et Greenleaf (1971).}

\pagerange{55}{57}
\subsection*{(2) Le découpage (pages 55 à 57)}

On découpe une surface conforme à bord $X$ le long d'un complexe topologique
$K$ de dimension $1$. La page demande quelle régularité imposer à $K$, et
propose la rectifiabilité.

\emph{Théorème 3.} Le découpé $\widetilde{X}$, muni du faisceau des fonctions
continues conformes dans $\operatorname{Int}\widetilde{X} \simeq X
\smallsetminus K$, est une surface conforme à bord.

La marge le qualifie de théorème d'uniformisation locale. Il revient à
trouver, autour de chaque point du bord de $\widetilde{X}$, un homéomorphisme
sur un demi-disque $\overline{\mathbb{H}} \cap \mathring{\mathbb{D}}$,
holomorphe à l'intérieur. Deux cas se présentent, selon que le point est au
milieu de $K$ (ordre au moins $2$) ou à une extrémité (ordre $1$) :

\emph{Lemme 2.} Soit $K \subset \mathbb{C}$ un segment topologique et $x \in
K$.
\begin{itemize}
  \item[a)] Si $x$ n'est pas une extrémité, il existe, pour chaque rive de
  $K$ en $x$, une application continue injective $f$ de
  $\overline{\mathbb{H}} \cap \mathring{\mathbb{D}}$ dans $\mathbb{C}$,
  holomorphe à l'intérieur, qui envoie le diamètre $]-1,1[$ dans $K$ et $0$
  sur $x$, et dont l'image est du côté de cette rive.
  \item[b)] Si $x$ est une extrémité, il existe $f$ continue sur
  $\overline{\mathbb{H}} \cap \mathring{\mathbb{D}}$, injective et holomorphe à
  l'intérieur, d'image dans $\mathbb{C} \smallsetminus K$, avec $f(0) = x$,
  et dont les restrictions à $]-1,0]$ et à $[0,1[$ sont injectives de même
  image contenue dans $K$.
\end{itemize}

Ces énoncés sont vrais pour tout arc de Jordan, sans hypothèse de
régularité.\footnote{C'est le théorème de Carathéodory (1913) : une
représentation conforme d'un domaine de Jordan se prolonge en un
homéomorphisme des adhérences, et plus généralement la représentation d'un
domaine simplement connexe à bord localement connexe se prolonge
continûment. Une fente de Jordan a deux rives, et son extrémité est atteinte
de façon unique. Le théorème 3 vaut donc pour tout complexe topologique
compact de dimension $1$ ; la question « rectifiabilité suffit ? » a la
réponse : elle n'est pas nécessaire. Là où la régularité compte, c'est au
recollement.}

\emph{Scholie.} Chaque rive de $K$ munit $K$ d'une structure analytique
réelle canonique, par les cartes au bord. Le passage de l'une à l'autre se
fait par l'identité, qui n'est pas analytique en général. Les propriétés de
régularité de cette application de recollement ($C^{\infty}$, analytique
réelle, etc.) reflètent celles de $K$ dans $\mathbb{C}$.

\pagerange{57}{60}
\subsection*{(3) Le recollement, et (4) les recollements bord à bord (pages 57 à 60)}

« On doit avoir un théorème dans le style du » :

\emph{Théorème 4} (énoncé attendu). Il s'agirait d'une équivalence entre deux
catégories. D'un côté, les triples $(X, K, S)$, où $X$ est une surface
conforme compacte à bord et $(K, S)$ un complexe topologique compact, de
sommets $S$, avec $\partial X \subset K \subset X$. De l'autre, les triples
$(\widetilde{X}, \widetilde{S}, \sigma)$, où $\widetilde{X}$ est une surface
conforme compacte à bord, $\widetilde{S}$ une partie finie de $\partial
\widetilde{X}$, et $\sigma$ une involution qui dit quels arcs de $\partial
\widetilde{X} \smallsetminus \widetilde{S}$ recoller à quels arcs. Les
composantes de $K \smallsetminus S$ peuvent être des cercles, précise la
marge. Le foncteur direct est le découpage (théorème 3). Le foncteur inverse
est connu topologiquement : $X = \widetilde{X}/\mathcal{R}$, avec pour
faisceau structural les fonctions continues conformes hors de $K$. Que cela
donne bien une structure conforme à bord revient à deux cas locaux (le
\emph{lemme 3}, en deux versions dont la première est barrée). Le premier
est le recollement de deux demi-disques le long d'un arc de leurs diamètres,
une couture qui passe par un point intérieur. Le second est le repliement
d'un demi-disque sur lui-même, les deux moitiés de son diamètre étant
identifiées par un homéomorphisme qui fixe $0$ : une couture qui s'arrête en
un point.\footnote{Les bornes de l'intervalle d'arrivée du cas b) sont
surchargées sur la page ; on donne le sens que la géométrie impose.}

Des deux côtés, la page demande que $K$ et $\sigma$ soient « rectifiables ».
Sous cette forme, l'équivalence n'est pas juste, et il faut dire ce qui
l'est.\footnote{Le recollement est ce qu'on appelle aujourd'hui la
\emph{soudure conforme} (\emph{conformal welding}). Pour un homéomorphisme
quasi-symétrique, la soudure existe et elle est unique ; la couture obtenue
est un quasi-arc, qui peut ne pas être rectifiable (Pfluger, 1960 ; Lehto et
Virtanen). Inversement, un arc rectifiable peut avoir une application de
recollement qui n'est pas quasi-symétrique : c'est le cas si l'arc a un
point de rebroussement. Les classes qui se correspondent sont celles des
quasi-arcs et des homéomorphismes quasi-symétriques, non celles des arcs et
homéomorphismes rectifiables. Pour des homéomorphismes plus généraux, la
soudure peut ne pas exister, ou ne pas être unique.} Le découpage n'exige
rien (théorème 3). Quand $K$ est rectifiable, $(X, K, S)$ est déterminé par
$(\widetilde{X}, \widetilde{S}, \sigma)$ : c'est l'unicité, et c'est là que
la rectifiabilité sert. L'existence du recollement pour une involution
$\sigma$ donnée est un problème de soudure conforme, qui se résout pour
$\sigma$ quasi-symétrique.

La note de marge datée « 12.1. », à la page 60, énonce l'outil de l'unicité :
si $C$ est une courbe dans un ouvert $U$ de $\mathbb{C}$ et $f$ une fonction
continue sur $U$, holomorphe dans $U \smallsetminus C$, alors $f$ est
holomorphe dans $U$.\footnote{La page porte « continue » là où il faut
« holomorphe dans $U \smallsetminus C$ » (lecture incertaine). L'énoncé est
vrai pour $C$ rectifiable, par le théorème de Morera (Painlevé) ; plus
généralement pour un ensemble de longueur $\sigma$-finie (Besicovitch,
1931). Il est faux pour une courbe de Jordan quelconque : il existe des
courbes de Jordan non effaçables, par exemple d'aire positive.} Appliquée à
l'application qui compare deux recollements, elle montre que ceux-ci
coïncident.

Le paragraphe (4), « recollements bord à bord de disques », est le cas du
théorème 4 où $K$ est une variété compacte de dimension $1$ et $S$ vide. On
y cherche des conditions naturelles sur $\sigma$. Pour une composante $C$ du
bord d'une surface conforme, deux données reviennent au même : une structure
de droite projective réelle sur $C$, et un « bouchage du trou » de bord $C$
par un disque. Car le groupe des automorphismes du disque agit sur son bord
comme le groupe projectif de $\mathbb{P}^{1}(\mathbb{R})$. Les propriétés de
régularité de cette structure projective reflètent celles de $C$ comme cercle
intérieur de la surface bouchée.\footnote{Ceci vaut sans réserve quand la
structure projective est analytique réelle pour la structure analytique de
$C$ : on recolle alors le long d'un difféomorphisme analytique, qui se
complexifie. Pour une régularité moindre, c'est encore la soudure conforme.}

\pagerange{60}{69}
\section*{VIII. Boucher les trous : disques, ombres, systèmes admissibles (pages 60 à 69)}

\pagerange{60}{62}
\subsection*{Un disque pointé, et la structure torique de son bord (pages 60 à 62)}

Une fois le trou bouché par un disque $D$ de bord $C$, choisir un point $a$ de
$\mathring D$ équivaut à choisir une \emph{structure torique} sur $C$
compatible à la structure de droite projective réelle. Le groupe structural
$\mathrm{PGL}_2(\mathbb{R})$ se réduit alors à son sous-groupe compact maximal
$\mathrm{O}(2)/\pm 1$. En effet, les points du disque correspondent aux
sous-groupes compacts maximaux du groupe des automorphismes, et le
stabilisateur de $a$ agit sur $C$ par rotations. On a alors l'inclusion
canonique $D \subset T_{X',a}$ de la section III.

\emph{Proposition.} Les deux catégories suivantes sont équivalentes :
\begin{itemize}
  \item[a)] les surfaces conformes compactes à bord, avec une composante $C$
  du bord munie d'une structure torique ;
  \item[b)] les surfaces conformes $X'$ munies d'un disque pointé $(D, a)$,
  $D \subset \operatorname{Int} X'$.
\end{itemize}
L'équivalence se fait en bouchant $C$ par le disque qui a la même structure
torique, et en retirant $\mathring D$ dans l'autre sens.\footnote{La page
demande que la structure torique soit « disons rectifiable » et que
$\partial D$ soit rectifiable dans $X'$. L'équivalence est acquise quand
la structure torique est analytique pour la structure analytique de $C$, et
$\partial D$ analytique dans $X'$ ; pour une régularité moindre, c'est la
question de soudure de la section VII.}

\pagerange{62}{63}
\subsection*{L'ombre d'un disque (pages 62 et 63)}

Pour $(X', a)$ fixé, chaque disque $D$ autour de $a$ définit dans
$T_{X',a}$ un disque euclidien, son \emph{ombre} $o(D)$. Les disques de
$T_{X',a}$ forment un torseur sous $\mathbb{R}^{*}_{+}$, un espace de dimension
$1$ ; l'application $D \mapsto o(D)$ n'a donc rien d'injectif. Elle est
compatible aux homothéties : les disques concentriques $\lambda D$,
$0 < \lambda \leqslant 1$, définis par l'inclusion $D \subset T_{X',a}$, ont
pour ombres les $\lambda\, o(D)$.

On prend alors $X'$ compacte, munie de points marqués $(a_i)_{i\in I}$ et
d'une famille de disques deux à deux disjoints $D_i$ autour des $a_i$, qui
ont pour ombres des $o(D_i) \subset T_{X',a_i}$. La question est de trouver
des conditions naturelles sur le système des $D_i$ qui impliquent qu'il soit
déterminé par le système des ombres.

\pagerange{64}{65}
\subsection*{Systèmes admissibles (pages 64 et 65)}

On laisse tomber le prime et l'on suppose $X$ connexe, $I$ non vide. La
surface $X^{*} = X \smallsetminus \bigcup_i \mathring{D}_i$ a une métrique
canonique à courbure constante (section IV). Dans le cas hyperbolique, qui
exclut $g = 0$ et $\nu = \operatorname{card} I \in \{1,2\}$, c'est la
métrique hyperbolique, et les bords sont géodésiques. Pour $(g,\nu) = (0,2)$,
$X^{*}$ est un anneau fermé, dont la métrique plate canonique rend encore les
bords géodésiques. Dans ces cas, chaque cercle $\partial D_i$ reçoit de la
métrique une structure torique canonique, par l'abscisse
curviligne.\footnote{La page dit « dans tous les cas ». Pour $(g,\nu) =
(0,1)$, $X^{*}$ est un disque fermé, dont les métriques distinguées dépendent
du choix d'un point (page 43, cas b')) : il n'y a pas de structure torique
canonique sur son bord. Une note de marge, peu lisible, mentionne aussi le
plan projectif réel « à un trou ».} Il est tentant d'exiger que les deux
structures toriques de chaque $\partial D_i$, celle qui vient de $D_i$ et celle
qui vient de $X^{*}$, coïncident. Cela signifie que $X$ se déduit de $X^{*}$ en
bouchant les trous par les disques canoniques. On dit alors que le système
$(D_i)$ est \emph{admissible} ; c'est la définition de la lettre à Bers.

Il pose deux questions sans y répondre. Comment l'admissibilité se lit-elle
dans la géométrie intrinsèque de $X$ ? Par exemple, quand $X$ a une métrique
canonique ($X \not\simeq \mathbb{P}^{1}$), quel rapport y a-t-il avec la
condition que les $D_i$ soient des disques géodésiques
$\{\mathrm{dist}(x, a_i) \leqslant r_i\}$ ?

\pagerange{65}{69}
\subsection*{L'espace de Teichmüller à bord, et une conjecture (pages 65 à 69)}

On fixe une surface topologique orientée de référence $X_0$, de type
$(g,\nu)$, et l'on considère l'espace de Teichmüller à bord
$\mathrm{TB}_{g,\nu}$. Il est homéomorphe à $\mathbb{R}^{6g-6+3\nu}$, ou
mieux à $\mathbb{C}^{3g-3+\nu} \times \mathbb{R}^{\nu}$. Le bouchage des
trous définit une application
\[
\mathrm{TB}_{g,\nu} \longrightarrow \mathrm{T}_{g,\nu}.
\]
Il présume qu'elle est une fibration de fibre $\mathbb{R}^{\nu}$. Pour une
sphère à trois trous, à combinatoire fixée, faire varier une, deux ou trois
des longueurs ne change pas la sphère bouchée à isomorphisme canonique près :
c'est $\mathbb{P}^{1}$ pointée en trois points. Il faudrait comprendre, dit-il,
comment les disques $D(l_i)$ qu'on retire dépendent des longueurs $l_i$, et en
particulier si chacun ne dépend que de sa propre longueur. C'est la question
que reprend le modèle de la lettre.

Admettant ce point, la famille des systèmes admissibles autour des $a_i$ d'une
surface pointée donnée est homéomorphe à $\mathbb{R}^{\nu}$, comme l'ensemble
des systèmes d'ombres possibles. D'où la

\emph{Conjecture} (page 67). Pour une surface conforme pointée, au moins dans
le cas hyperbolique, l'application $(D_i) \mapsto (o(D_i))$, des systèmes
admissibles de disques autour des $a_i$ vers les systèmes de disques
euclidiens des $T_{X,a_i}$, est bijective.

Elle ferait de $\mathrm{TB}_{g,I} \to \mathrm{T}_{g,I}$ un fibré principal
de groupe $(\mathbb{R}^{*}_{+})^{I}$.\footnote{Sous cette forme, la
conjecture ne peut pas tenir, pour une raison élémentaire ; l'argument est le
nôtre. Sur $\mathbb{P}^{1}$ pointée en $0, 1, \infty$, un disque $D_0$ autour
de $0$, disjoint des disques autour de $1$ et de $\infty$, ne contient ni $1$
ni $\infty$. Le théorème du quart de Koebe borne alors son rayon conforme en
$0$, mesuré dans la coordonnée $z$, par $4\,\mathrm{dist}(0, \mathbb{C}
\smallsetminus D_0) \leqslant 4$. Les ombres restent donc dans une partie
bornée de $T_{X,0}$, et l'application n'est pas surjective. En passant au
revêtement universel (lemme de Schwarz en genre au moins $2$, Koebe en genre
$1$), on borne de même les ombres pour tout type. La lettre à Bers du 15 avril
(pages 17 et 18) enregistre le phénomène en remplaçant la bijectivité par des
inégalités $r_i < \rho_i$.} Si la relation présumée plus haut était vraie, ce
fibré aurait une section canonique : des métriques euclidiennes canoniques sur
les $T_{X,a_i}$. En genre au moins $1$, cela ne l'inquiète pas, puisque $X$ a
une métrique canonique. En genre $0$, cela lui semble étrange, puisque
$\mathbb{P}^{1}$ n'en a pas.

\pagerange{69}{73}
\section*{IX. Retour sur le principe de chirurgie : le lemme 4 (pages 69 à 73)}

Le titre, « 5. Retour sur le principe de chirurgie conforme », est daté
« 12.1. » en marge. Les lemmes locaux 2 et 3, « ingrédient technique
essentiel pour le théorème 4 (et lui étant équivalents) », sont formulés de
façon ambiguë. Il vaut mieux se ramener à des énoncés globaux
d'uniformisation par le disque $\mathbb{D}$. Soit $\zeta = e^{2i\pi/3}$. Les
points $1, \zeta, \bar\zeta$ découpent le cercle $\partial\mathbb{D}$ en trois
arcs : $\alpha$, de $\zeta$ à $\bar\zeta$ sans passer par $1$ ; $\beta$, de
$1$ à $\zeta$ ; $\bar\beta$, de $1$ à $\bar\zeta$.

\emph{Lemme 4.} A) Soit $V$ le domaine intérieur d'une courbe de Jordan $C$,
et $I$ un arc topologique qui rencontre $C$ en une seule de ses extrémités
$b$, son autre extrémité $a$ étant dans $V$. Soit $U = V \smallsetminus I$.
Il existe une application continue $f : \mathbb{D} \to \overline{U}$ telle
que :
\begin{itemize}
  \item[1)] $f$ induit un homéomorphisme $\mathring{\mathbb{D}} \simeq U$, et
  elle est holomorphe sur $\mathring{\mathbb{D}}$ ;
  \item[2)] $f(\zeta) = f(\bar\zeta) = b$, et $f$ induit une bijection
  $\alpha \smallsetminus \{\zeta,\bar\zeta\} \simeq C \smallsetminus \{b\}$ ;
  \item[3)] $f(1) = a$, et $f$ induit des bijections de $\beta$ et de
  $\bar\beta$ sur $I$.
\end{itemize}
Autrement dit, $f$ est l'uniformisation de $\overline{V}$ découpé le long de
$I$. Elle est unique, parce que le groupe des automorphismes de $\mathbb{D}$
agit de façon simplement transitive sur les triples de points distincts du
cercle rangés dans l'ordre de l'orientation. Il en résulte une bijection
canonique
\[
\varphi : \beta \xrightarrow{\ \sim\ } \bar\beta, \qquad \varphi(1) = 1,
\quad \varphi(\zeta) = \bar\zeta,
\]
l'application de recollement des deux rives de $I$.\footnote{La page suppose
$C$ et $I$ rectifiables. Ce n'est pas nécessaire pour A) : c'est le théorème
de représentation conforme avec le théorème de Carathéodory, puisque le bord
de $U$ est localement connexe. La page appelle les conditions tantôt 1) à 4),
tantôt a), b), c).}

B) Si $C$ et $I$ sont rectifiables, $\varphi$ et son inverse sont absolument
continus.\footnote{La page dit « $\varphi$ est rectifiable » sans définir
ce mot pour un homéomorphisme. La lecture « absolument continu ainsi que
son inverse » rend l'assertion vraie, par le théorème de F. et M. Riesz :
sur un bord rectifiable, la mesure harmonique et la longueur sont
mutuellement absolument continues.} Réciproquement, la page affirme que tout
homéomorphisme $\varphi : \beta \to \bar\beta$ de cette classe provient d'une
telle situation. On normalise $\overline{V}$ en $\mathbb{D}$, puis $a$ en $0$
et $b$ en $1$. On cherche alors une application continue $f_\varphi :
\mathbb{D} \to \mathbb{D}$, injective et holomorphe sur $\mathring{\mathbb{D}}$,
avec $f(1) = 0$ et $f(\zeta) = f(\bar\zeta) = 1$. Elle doit induire un
homéomorphisme de $\alpha \smallsetminus \{\zeta,\bar\zeta\}$ sur
$\partial\mathbb{D} \smallsetminus \{1\}$, envoyer $\beta$ et $\bar\beta$ sur
un même arc $I$ joignant $0$ à $1$, avec $I \smallsetminus \{1\} \subset
\mathring{\mathbb{D}}$, et vérifier $\varphi = (f|_{\bar\beta})^{-1} \circ
f|_{\beta}$.\footnote{La page écrit $f(0) = 0$. C'est impossible, puisque $0$
est l'extrémité $a$ de la fente et n'est pas dans l'image de
$\mathring{\mathbb{D}}$ ; comme en A), c'est $f(1) = a = 0$.} Cette
réciproque est un problème de soudure conforme, et la page ne la démontre
pas. On sait la résoudre pour $\varphi$ quasi-symétrique : la fente obtenue
est alors un quasi-arc, qui peut ne pas être rectifiable. Les homéomorphismes
quasi-symétriques ne sont d'ailleurs pas tous absolument continus.

\emph{Scholie.} La page met en correspondance biunivoque quatre données :
\begin{itemize}
  \item[a)] les homéomorphismes rectifiables $\psi : [0,1] \to [0,1]$ avec
  $\psi(0) = 0$ ;
  \item[b)] les arcs rectifiables $I \subset \mathbb{D}$ d'extrémités $0$ et
  $1$, avec $I \cap \partial\mathbb{D} = \{1\}$ ;
  \item[c)] les applications continues $f : \mathbb{D} \to \mathbb{D}$ qui ont
  les propriétés de B) ;
  \item[d)] les séries entières $\sum_{i \geqslant 0} a_i z^{i}$, de rayon de
  convergence au moins $1$ et de module au plus $1$ sur
  $\mathring{\mathbb{D}}$, qui se prolongent par continuité en une telle
  application.
\end{itemize}
Les correspondances b) $\leftrightarrow$ c) $\leftrightarrow$ d) sont
acquises : représentation conforme de $\mathbb{D} \smallsetminus I$
normalisée, et théorème de Carathéodory. L'application b) $\to$ a), qui
associe à une fente son recollement, est injective, parce qu'un arc
rectifiable est effaçable (section VII). Sa surjectivité sur une classe
d'homéomorphismes est la soudure conforme, et « rectifiable » n'est pas la
classe pour laquelle on sait répondre.

\pagerange{73}{81}
\section*{X. Des surfaces à bord aux courbes stables (pages 73 à 81)}

Le paragraphe 6 est daté « 13.1. ». Soit $X$ une surface holomorphe compacte à
bord, connexe, de type $(g,\nu)$, et $C$ un système de courbes le long
desquelles la découper : un objet du groupoïde de Teichmüller à bord. On
découpe $X$ le long de $C$ et l'on obtient $X'$. On en bouche canoniquement
les trous, grâce aux structures toriques que la métrique canonique de $X'$
donne à son bord. On obtient une surface holomorphe compacte sans bord $X''$,
munie de disques pointés. Ces disques se groupent par paires, sauf les $\nu$
disques qui viennent de $\partial X$, et chaque paire est munie de
l'anti-isomorphisme unitaire entre ses bords qui recollait les deux rives.

On voit alors $X''$ comme la normalisée d'une courbe algébrique complexe
\emph{stable} $X'''$, au sens de Deligne et Mumford. On identifie les centres
des disques appariés, et l'on marque les centres des autres. En chaque point
double $x$ de $X'''$, qui correspond à une paire $\{x_1, x_2\}$ de points de
$X''$, l'anti-isomorphisme des disques donne un élément non nul
\[
t_x \in T_{X'',x_1} \otimes_{\mathbb{C}} T_{X'',x_2},
\qquad\text{soit}\qquad
t_x \in T^{*}_{X'',x_1} \wedge_{\mathbb{C}^{*}} T^{*}_{X'',x_2},
\]
où $T^{*}$ désigne le $\mathbb{C}^{*}$-torseur des vecteurs non nuls. C'est
ce qu'il appelle une \emph{structure de propagation normale}. La
correspondance est celle de la page 52 : la donnée d'une paire de disques et
d'un anti-isomorphisme de leurs bords équivaut à celle de l'un des deux
disques et d'un élément de $T_{x_1} \otimes T_{x_2}$.\footnote{Recoller les
voisinages de $x_1$ et $x_2$ par $zw = t$ est la construction qu'on appelle
aujourd'hui \emph{plombage} (\emph{plumbing}), et $t$ est le paramètre de
lissage du nœud. La page ne le nomme pas.}

On obtient ainsi un foncteur. Il va des surfaces holomorphes compactes à bord
de type $(g,\nu)$, munies de leurs disques, vers les courbes algébriques
complexes stables de type $(g,\nu)$, munies d'une structure de propagation
normale en leurs points doubles et d'un système admissible de disques autour
de leurs points marqués.\footnote{Un mot en interligne, lu « quasi-vides »
sous toute réserve, qualifie la première catégorie ; on ne l'utilise pas. Les
flèches de la première catégorie sont les isomorphismes.} Au niveau des
multiplicités modulaires, analytiques réelles, on a sûrement un morphisme
\[
\mathrm{MDB}_{g,\nu} \longrightarrow \mathrm{M_{st}ND}_{g,\nu}.
\]
La source classe les surfaces holomorphes compactes « disquées » à bord de
type $(g,\nu)$, le but les courbes stables à propagation normale et à
système admissible de disques. Le morphisme est compatible aux flèches vers la
catégorie $\mathrm{mdt}_{g,\nu}$ des types combinatoires : les diagrammes de
Mumford--Deligne, que l'on appelle aujourd'hui graphes duaux, pondérés par les
genres, avec des structures toriques aux arêtes (d'où le $t$).

\begin{tikzcd}
\mathrm{MDB}_{g,\nu} \arrow[rr] \arrow[dr] & & \mathrm{M_{st}ND}_{g,\nu} \arrow[dl] \\
& \mathrm{mdt}_{g,\nu} &
\end{tikzcd}

De même, sans bord, ou avec un bord générique réduit à $\nu$ points, et sans
disques, on aurait un morphisme de multiplicités algébriques complexes
$\mathrm{MD}_{g,\nu} \to \mathrm{M_{st}N}_{g,\nu}$ au-dessus de
$\mathrm{md}_{g,\nu}$, par le même triangle.

Les dimensions se correspondent : réelle, celle de
$\mathbb{C}^{3g-3+\nu} \times \mathbb{R}^{\nu}$ dans le premier cas ;
complexe, $3g - 3 + \nu$ dans le second. Une courbe stable à $k$ points
doubles varie en effet avec $3g - 3 + \nu - k$ paramètres complexes, auxquels
les $t_x$ ajoutent $k$ paramètres complexes et les disques $\nu$ paramètres
réels.

\emph{Question} (page 77). Ces flèches sont-elles des isomorphismes, et
comment expliciter un inverse ? Si la détermination conjecturale des systèmes
admissibles était correcte, on construirait aisément la flèche inverse. Un
$u_i \in T^{*}_{x'_i} \wedge_{\mathbb{C}^{*}} T^{*}_{x''_i}$ équivaut à un
accouplement dualisant $w_i : T_{x'_i} \otimes T_{x''_i} \to \mathbb{C}$. On
choisit un disque $D'_i$ dans $T_{x'_i}$, d'où le disque polaire $D''_i$ dans
$T_{x''_i}$ ; on réalise ces disques dans la normalisée $X''$, on retire les
disques ouverts, et l'on recolle les bords par l'anti-isomorphisme que définit
$w_i$. On obtient une surface à bord $X$. Les coutures $\Gamma_i$ ne sont peut-être
pas géodésiques --- « à examiner de plus près » --- mais elles définissent des
géodésiques uniques, par les principes de Thurston.

Il faudrait voir que le résultat ne dépend pas du choix des paires
mutuellement polaires. L'argument heuristique est de remplacer $D'_i$ par
$\lambda D'_i$ et $D''_i$ par $\lambda^{-1} D''_i$, avec $\lambda$ voisin de
$1$ : on obtiendrait la même surface $X$, les $\Gamma_i$ étant remplacés par
des cercles voisins. Il ajoute : « c'est probablement tout à fait faux
! ».\footnote{Ce qui est sûr est la partie qui concerne la surface. Des
disques concentriques sont des disques de coordonnées d'une même carte, et le
recollement par $zw = t$ de leurs complémentaires ne dépend pas des rayons,
tant que les disques restent dans le domaine des cartes : c'est
l'indépendance, classique, du plombage par rapport aux rayons. Le doute de la
page porte sans doute sur ce que deviennent les disques et leurs coutures,
qu'on ne sait pas contrôler faute de connaître les systèmes admissibles.}

Le cas-clef est $(g,\nu) = (0,4)$ : deux pantalons de Thurston recollés.
D'un côté, les pantalons de types combinatoires $\Sigma' \sqcup \{s\}$ et
$\Sigma'' \sqcup \{t\}$, recollés suivant $\{s,t\}$, avec les longueurs des
quatre bords et de la couture et un angle de torsion $\theta$ : six
paramètres réels. De l'autre, les paires de pantalons généralisés dont les
bords $s$ et $t$ sont réduits à des points, avec leurs quatre périmètres
$\lambda_i$ et une identification $\mathbb{C}^{*}$ de paramètre de lissage :
six encore.\footnote{L'indice de la couture et le groupe
$\mathbb{U}_{\Sigma' \wedge \Sigma''}$ où vit $\theta$ sont lus sous
réserve. En termes actuels, c'est la comparaison entre les coordonnées de
Fenchel--Nielsen (longueur et torsion de la couture) et les coordonnées de
plombage (le paramètre $t$).} Une fois ce cas compris, il faudrait, pour un
type de diagramme donné, comprendre l'application induite sur les groupes
fondamentaux par $\mathrm{MDB} \to \mathrm{M_{st}N}_{g,\nu}$. Celui du but
s'exprime par des groupes de Teichmüller spéciaux, et il « doit en être de
même pour le premier ». La page 81, qu'il numérote 15, achève la dernière
phrase : il ne devrait pas être bien difficile de vérifier que cette
application est étale, en calculant ou en interprétant l'application
tangente.\footnote{« Étale » est une lecture incertaine, et le raccord entre
la fin de la page 80 et les deux lignes de la page 81 est probable, non
certain.}

\pagerange{83}{88}
\section*{XI. Modules $M_{0,4}$ et l'octaèdre (pages 83 à 88)}

Soit $k$ un corps algébriquement clos de caractéristique différente de $2$,
$I$ un ensemble à quatre éléments et $J = \mathfrak{P}_{2,2}(I)$ l'ensemble,
à trois éléments, de ses partitions en deux paires. Le groupe de Klein
$V(J) \subset \mathfrak{S}_I$, formé de l'identité et des trois doubles
transpositions, est un plan vectoriel sur $\mathbb{F}_2$ dont les trois
éléments non nuls correspondent aux éléments de $J$.\footnote{La page écrit
$V(J)^{*} \simeq J$ ; l'astérisque désigne ici l'ensemble des éléments non
nuls, comme le $V^{\times}$ de la page 85, et non le dual.} On s'intéresse
aux droites projectives $X$ sur $k$ munies d'un plongement
$I \hookrightarrow X$.

L'action de $V(J)$ sur $I$ se prolonge en une action sur $X$ : chaque double
transposition de quatre points distincts est induite par une unique
homographie, puisque le birapport est invariant par le groupe de Klein. On a
ainsi $V(J) \hookrightarrow G_X = \mathrm{Aut}(X)$. Chaque involution non
triviale a deux points fixes, et les trois paires sont disjointes : les points
fixes forment un ensemble $S$ de six points. Toutes les actions fidèles d'un
plan vectoriel sur $X$ sont conjuguées, même à base fixée $(e_0, e_1)$, au
modèle
\[
e_0 \mapsto -z, \qquad e_1 \mapsto 1/z, \qquad e_0 + e_1 \mapsto -1/z,
\]
de points fixes $\{0,\infty\}$, $\{\pm 1\}$, $\{\pm i\}$. Les sous-groupes de
Klein de $G_X$ correspondent donc aux parties de six points isomorphes à
$\bigl(\mathbb{P}^{1}, \{0,\infty,\pm 1,\pm i\}\bigr)$, les sommets d'un
octaèdre. Le normalisateur $\Gamma$ d'un tel $V$ est isomorphe à
$\mathfrak{S}_4$ : c'est le groupe octaédral. Si $V = V(J)$, on a un
isomorphisme d'extensions

\begin{tikzcd}[column sep=small]
1 \arrow[r] & V \arrow[r] \arrow[d, no head, "\Vert" description] & \mathfrak{S}_{I} \arrow[r] \arrow[d, "\wr"] & \operatorname{Aut}(V) \arrow[r] \arrow[d, no head, "\Vert" description] & 1 \\
1 \arrow[r] & V \arrow[r] & \Gamma \arrow[r] & \operatorname{Aut}(V) \arrow[r] & 1
\end{tikzcd}

défini à un automorphisme intérieur par un élément de $V$ près.\footnote{Ces
isomorphismes forment un torseur sous $V$, parce que $H^{1}(\mathfrak{S}_3,
V) = 0$ : la restriction à un $2$-Sylow, sur lequel $V$ est libre, est
injective et nulle. La vérification est la nôtre.}

Les scindages de l'extension $1 \to V \to \Gamma \to \operatorname{Aut}(V) \to
1$ correspondent aux quatre paires de faces antipodales de l'octaèdre. Les
faces sont les parties $S' \subset S$ à trois éléments telles que
$S' \cup \mathsf{a}(S') = S$, où $\mathsf{a}$ est l'antipodie qui
échange les deux points fixes de chaque involution de $V$. Le stabilisateur
dans $\Gamma$ d'une paire $\{S', \mathsf{a}(S')\}$ est un groupe diédral
d'ordre $6$, supplémentaire de $V$.\footnote{La page dit « le sous-groupe de
$\Gamma$ qui l'invarie », en parlant d'une face $S'$. Le stabilisateur d'une
face seule est cyclique d'ordre $3$ ; c'est celui de la paire de faces
antipodales qui est un scindage.} Le groupe $\Gamma$ est le groupe des
automorphismes de l'octaèdre combinatoire qui conservent l'orientation ;
choisir une orientation revient à choisir une racine carrée de $-1$, en
déclarant directe la face $(0, 1, \sqrt{-1})$. Notant
$\operatorname{Oct}(V)$ l'octaèdre et $I(V)$ l'ensemble de ses quatre paires
de faces antipodales, on a $\Gamma(V) \simeq \mathfrak{S}_{I(V)}$, et $V$ agit
simplement transitivement sur $I(V)$, qui devient un $V$-torseur.

\emph{Proposition} (page 85). Soit $k$ un corps de caractéristique
différente de $2$ où tout élément est un carré, $X$ une droite projective sur
$k$ et $G = \mathrm{Aut}_k X$. Il y a des correspondances biunivoques entre :
\begin{itemize}
  \item[a)] les sous-groupes $V$ de $G$ isomorphes au groupe de Klein ;
  \item[b)] les sous-groupes $\Gamma$ de $G$ isomorphes à $\mathfrak{S}_4$ ;
  \item[c)] les parties octaédrales $S \subset X(k)$, isomorphes à
  $\bigl(\mathbb{P}^{1}_k, \{0,\infty,\pm 1,\pm i\}\bigr)$.
\end{itemize}
On passe de l'une à l'autre par $\Gamma = \operatorname{Nor}_G(V)$ ; par
« $V$ = l'unique sous-groupe distingué de $\Gamma$ isomorphe au groupe de
Klein » ; et par « $S$ = l'ensemble des points fixes des éléments non
triviaux de $V$ ».\footnote{L'hypothèse sur le corps est en partie
illisible ; la page parle de racines carrées et de caractéristique
différente de $2$, et c'est ce qu'il faut pour que les points fixes des
involutions et les éléments du normalisateur soient rationnels. Le mot lu
« réelle » devant « droite projective » ne peut pas être gardé : sur
$\mathbb{R}$, $\mathrm{PGL}_2(\mathbb{R})$ ne contient pas de
$\mathfrak{S}_4$, et le groupe $\{z, -z, 1/z, -1/z\}$ a deux points fixes non
réels.} L'ensemble des points fixes d'éléments non triviaux de $\Gamma$ est
plus grand : outre les six sommets, il comprend les douze milieux d'arêtes et
les huit centres de faces, les orbites exceptionnelles du groupe octaédral.
La page 86 caractérise les structures octaédrales sur un ensemble $S$ de six
points muni d'une action de $V$. Pour chaque $\sigma \in V$ non trivial,
$S^{\sigma}$ doit avoir deux éléments, $S$ doit être la réunion, forcément
disjointe, des $S^{\sigma}$, et chaque $\sigma$ doit échanger les deux points
de chacune des deux autres paires. Cela revient à munir $S$ d'une involution
sans point fixe dont les $S^{\sigma}$ sont les orbites.

Une action fidèle de $V$ sur $X$ étant donnée, $V$ agit librement hors de $S$,
et chaque orbite $I$ de $V$ dans $X \smallsetminus S$ est un $V$-torseur, à ne
pas confondre avec le torseur $I(V)$ des paires de faces. Soit $X(I)$ une
droite projective de référence munie d'une action fidèle de $V(I)$, que la
page construit à partir de $I$.\footnote{La construction, à la page 86, est
presque entièrement raturée, et l'on ne sait pas si elle dépend du choix
d'une racine carrée de $-1$. Les énoncés qui suivent valent pour tout choix
d'une telle droite de référence, le torseur $\mathsf{t}$ dépendant du
choix.} Pour deux $V$-torseurs, on a $X(T \wedge T') = X(T) \wedge^{V} T'$ (produit
contracté). Si $\mathsf{t}$ est le $V$-torseur tel que $I(V) \simeq I
\wedge^{V} \mathsf{t}$, on obtient
\[
X \simeq X(I(V)) = X(I) \wedge^{V} \mathsf{t}, \qquad
X/V \simeq X(I)/V \simeq \Sigma(J),
\]
où $\Sigma(J)$ est une droite projective qui ne dépend que de
$J$.\footnote{La lettre de $\Sigma(J)$ est tracée comme un $X$ barré ; la
transcription la lit $\Sigma$. Le quotient d'une droite projective par un
groupe de Klein est une droite projective, et le revêtement
$X \to X/V$, de degré $4$, est ramifié au-dessus de trois points,
images des trois paires $S^{\sigma}$, indexés par $J$.} Pour $I$ donné, se
donner $X$ et un plongement $V \hookrightarrow G_X$ revient à se donner un
$V$-torseur $\mathsf{t}$, par $\mathsf{t} \mapsto X(I) \wedge^{V}
\mathsf{t}$. Se donner $(X, I \hookrightarrow X)$ revient à se donner un
point $\xi$ de $\Sigma(J)^{*}$, le complémentaire des trois points de
ramification, et le torseur $\mathsf{t} = I \wedge X(I)_\xi^{-1}$, de
sorte que la fibre $X_\xi = X(I)_\xi \wedge \mathsf{t}$ soit $I$.

\emph{Théorème} (page 88). Soit $k$ algébriquement clos de caractéristique
différente de $2$. Le groupoïde des couples $(X, I \hookrightarrow X)$, où
$X$ est une droite projective sur $k$ et $I$ un ensemble à quatre éléments,
est équivalent à celui des couples $(I, \xi)$ avec $\xi \in
\Sigma(J(I))^{*}$. L'équivalence est donnée par
\[
X = X(I) \wedge^{V(I)} t(\xi, I), \qquad t(\xi, I) = I \wedge
X(I)_\xi^{-1},
\]
et $I \hookrightarrow X$ est la fibre de $X \to X/V(I) \simeq \Sigma(J)$
au-dessus de $\xi$.

C'est la description de l'espace des modules $M_{0,4}$ : une droite
projective canoniquement attachée à l'ensemble $J$ des trois partitions de
$I$, privée de trois points indexés par $J$. Ces trois points sont les trois
dégénérescences stables d'une droite à quatre points, et le groupe
$\mathfrak{S}_4$ agit à travers $\mathfrak{S}_3 = \mathfrak{S}_J$, le groupe
de Klein agissant trivialement.\footnote{Les automorphismes se correspondent
bien : un couple $(X, I \hookrightarrow X)$ générique a pour automorphismes
les éléments de $V$, réalisés par des homographies qui permutent $I$, et un
couple $(I, \xi)$ générique aussi, puisque $V$ agit trivialement sur
$\Sigma(J)$. La vérification est la nôtre.} Si $I$ a un point marqué, $X$
s'identifie à $X(I)$ dès qu'on choisit un point de $X(I)$ au-dessus de $\xi$ ;
plus généralement, un isomorphisme « admissible » $X \simeq X(I)$ revient à un
$V$-isomorphisme $I \simeq X(I)_\xi$.

\pagerange{90}{100}
\section*{XII. Courbes elliptiques réelles (pages 90 à 100)}

La page 90 est datée « Déc 83 ».

\pagerange{90}{91}
\subsection*{Modules munis d'une involution (pages 90, 91 et 99)}

Soit $k$ un anneau commutatif, et $M$ un $k$-module muni d'un automorphisme
$\sigma$ tel que $\sigma^{2} = \mathrm{id}$, avec $M$ \emph{$2$-régulier} :
la multiplication par $2$ y est injective. On pose
\[
P = M_{+} = \operatorname{Ker}(1 - \sigma), \qquad Q = M_{-} =
\operatorname{Ker}(1 + \sigma).
\]
On a $P \cap Q = 0$, par $2$-régularité, et $2x = (x + \sigma x) + (x -
\sigma x)$, d'où
\[
2P \oplus 2Q \subset M' = 2M \subset P \oplus Q.
\]
Donc $M'$ est l'image réciproque d'un sous-$k/2k$-module $m \subset P/2P
\oplus Q/2Q$. De $2x \in P \Rightarrow x \in P$, on tire $M' \cap P = 2P$, et
de même $M' \cap Q = 2Q$ ; ainsi $m$ rencontre trivialement $p = P/2P$ et
$q = Q/2Q$. C'est donc le graphe d'un isomorphisme entre un sous-module de
$p$ et un sous-module de $q$. Réciproquement, $(M, \sigma)$ se reconstitue à
partir de $(P, Q, m)$ : $M'$ est l'image réciproque de $m$, $\sigma$ y est
induite par $\mathrm{id}_P \oplus (-\mathrm{id}_Q)$, et la multiplication par
$2$ est un isomorphisme de $(M,\sigma)$ sur $(M', \sigma)$.

\emph{Proposition.} Le foncteur $(M, \sigma) \mapsto (P, Q, m)$ est une
équivalence entre la catégorie des modules $2$-réguliers à involution et celle
des triples formés de deux modules $2$-réguliers et d'un sous-module $m$ de
$P/2P \oplus Q/2Q$ qui rencontre trivialement chacun des deux
facteurs.\footnote{La page 91 est le brouillon de ce calcul, sur un listing
daté du 19 août 1982. Elle écrit $\sigma(x,y) = x - y$ pour
$\sigma(x,y) = (x, -y)$, et une partie en est barrée.}

\emph{Corollaire} (page 99). Supposons $k$ principal, $2$-régulier, avec
$k_0 = k/2k$ un corps, et $M$ libre de rang $2$ avec $\sigma \neq \pm 1$.
Alors $P$ et $Q$ sont libres de rang $1$, et $m$ est soit nul, soit une
droite de $p \oplus q$ distincte de $p$ et de $q$. Si $k^{*} \to k_0^{*}$ est
surjectif, on peut la ramener à la diagonale. Il y a donc exactement deux
classes de conjugaison d'éléments d'ordre $2$ de $\mathrm{GL}_2(k)$ distincts
de $\pm 1$ :
\[
\begin{pmatrix} 1 & 0 \\ 0 & -1 \end{pmatrix} \quad (m = 0),
\qquad
\begin{pmatrix} 1 & 1 \\ 0 & -1 \end{pmatrix} \quad (m \text{ diagonale}).
\]
La seconde s'obtient avec $M' = \{(x,y) \in k^{2} \mid x - y \in 2k\}$, dans
la base $u_0 = (2, 0)$, $v_1 = (1,1)$, où $\sigma u_0 = u_0$ et $\sigma v_1 =
u_0 - v_1$. Pour $k = \mathbb{Z}$, c'est la réflexion d'un réseau
hexagonal par rapport à la droite d'un sommet $u$, dans la base de deux
sommets consécutifs $u, v$. Plus généralement, si $u_{+}$ et $u_{-}$ sont des
bases de $\Pi_{+}$ et $\Pi_{-}$, le réseau $\Pi$ est $\mathbb{Z}u_{+} +
\mathbb{Z}u_{-}$ dans le premier cas ; dans le second, c'est l'ensemble des
$x u_{+} + y u_{-}$ avec $x, y \in \frac12\mathbb{Z}$ et $x - y \in \mathbb{Z}$.

\pagerange{92}{92}
\subsection*{Deux familles (page 92)}

Une courbe elliptique réelle équivaut à un système $(\Pi, \sigma, c)$ :
$\Pi$ est un $\mathbb{Z}$-module libre de rang $2$, $c$ une structure
conforme orientée sur $\Pi_{\mathbb{R}} = \Pi \otimes \mathbb{R}$, et $\sigma$
une involution de $\Pi$ qui conserve la structure conforme et renverse
l'orientation. En effet, la structure conforme orientée fait de
$\Pi_{\mathbb{R}}$ une droite complexe, la courbe est
$\Pi_{\mathbb{R}}/\Pi$, et $\sigma$ en est la conjugaison. Pour le couple
$(\Pi, \sigma)$, il y a deux possibilités à isomorphisme près : les deux
classes du corollaire, qui sont isogènes. Une structure conforme invariante
par $\sigma$ rend orthogonales les droites propres $\Pi_{+} \otimes
\mathbb{R}$ et $\Pi_{-} \otimes \mathbb{R}$. On la normalise par
$Q(e_{+}) = 1$, où $e_{\pm}$ engendre $\Pi_{\pm}$ :
\[
Q(x e_{+} + y e_{-}) = x^{2} + a y^{2}, \qquad a \in \mathbb{R}^{*}_{+}.
\]
D'où deux familles de courbes elliptiques réelles, $E_a$ (où $\Pi =
\Pi_{+} \oplus \Pi_{-}$) et $E'_a$ (où $\Pi_{+} \oplus \Pi_{-}$ est d'indice
$2$ dans $\Pi$). Les automorphismes d'une courbe elliptique réelle sont
$\pm 1$.\footnote{La seconde moitié de la page est très raturée, et c'est la
conclusion qui s'y lit (« l'identité, et $\sigma$ », puis
« l'automorphisme universel $-1$ »). Elle est juste : $\sigma$ agit sur
$\mathrm{Aut}(E_{\mathbb{C}}) = \mu_n$ par $\zeta \mapsto \zeta^{-1}$, et
seuls $\pm 1$ commutent à $\sigma$.}

\pagerange{93}{97}
\subsection*{Les structures réelles d'une courbe complexe donnée (pages 93 à 97)}

Soit $X = V/\Pi$ une courbe elliptique complexe. Ses structures réelles, en
tant que courbe elliptique, correspondent aux droites réelles de $V$ telles
que la réflexion orthogonale par rapport à elles conserve $\Pi$. Soit $G$ le
groupe des isométries de $\Pi$ pour la structure conforme, et $G^{+} =
\mathrm{Aut}(X)$, cyclique d'ordre $n \in \{2,4,6\}$. $X$ admet une
structure réelle si et seulement si $\det : G \to \{\pm 1\}$ est
surjectif.\footnote{C'est-à-dire si et seulement si l'invariant $j(X)$ est
réel ; la page ne le dit pas sous cette forme.} Alors $G$ est diédral
d'ordre $2n$, et les structures réelles sont les éléments de $G^{-}$, qui est
un torseur sous $G^{+}$. Deux structures réelles sont isomorphes si elles
sont conjuguées par $G^{+}$, qui agit à travers $G^{+}/\pm 1$. Il y a
toujours exactement deux orbites.
\begin{itemize}
  \item Pour $n = 2$ (cas non harmonique), $G^{+} = \{\pm 1\}$ est central
  et les deux orbites sont $\{\sigma\}$ et $\{-\sigma\}$.
  \item Pour $n = 4$ et $n = 6$, ce sont les réflexions par rapport aux
  sommets et par rapport aux côtés du polygone fondamental, carré ou
  hexagone.
\end{itemize}

\emph{Proposition} (page 94). Une courbe elliptique complexe à invariant $j$
réel a exactement deux structures réelles non isomorphes.\footnote{C'est le
classique : les deux formes réelles sont une courbe et sa tordue
quadratique par $-1$. La classification des courbes elliptiques réelles est
exposée par N. Alling, \emph{Real elliptic curves} (1981). La page ne la
cite pas.}

La page 94 présume ensuite que l'une des deux orbites est formée de
structures $E_a$ et l'autre de structures $E'_a$, avec le même $a$. La
suite du calcul la contredit. Le passage de $\sigma$ à $-\sigma$ échange
$\Pi_{+}$ et $\Pi_{-}$, conserve le type ($E$ ou $E'$), et donne
\[
a(-\sigma) = \frac{1}{a(\sigma)}, \qquad a(\sigma) =
\frac{Q(e_{-})}{Q(e_{+})},
\]
de sorte que $a(-\sigma) = a(\sigma)$ si et seulement si $a = 1$, le cas
carré. Dans le cas carré ($n = 4$), avec $e_0, e_1$ orthonormés dirigés vers
deux sommets, la réflexion par rapport à $\mathbb{R}e_0$ est du type $E_1$
(« sommets »). Celle qui échange $e_0$ et $e_1$ a pour invariants et
anti-invariants $e_0 + e_1$ et $e_0 - e_1$, qui engendrent un sous-réseau
d'indice $2$ : elle est du type $E'_1$ (« côtés »). Dans le cas hexagonal
($n = 6$), avec $u, v$ deux sommets consécutifs, la réflexion par rapport à
$\mathbb{R}u$ a pour invariants $\mathbb{Z}u$ et pour anti-invariants
$\mathbb{Z}(2v - u)$, d'indice $2$ : elle est du type $E'_a$ avec
\[
a = \frac{Q(2v - u)}{Q(u)} = 3,
\]
et $-\sigma$ donne $E'_{1/3}$.\footnote{La page écrit $a = \sqrt{3}$, d'où
$E'_{\sqrt{3}}$ et $E'_{1/\sqrt3}$. Mais $Q$ est une forme quadratique,
$Q(u) = Q(v) = 1$ et $u\cdot v = \frac12$ donnent $Q(2v - u) = 4 - 2 + 1 =
3$ ; c'est d'ailleurs $a = 3$ que demande la formule de la page 98,
$\tau = \frac12(1 + \sqrt{a}\, i) = e^{i\pi/3}$.}

\emph{En résumé} (page 97). Sauf dans le cas carré, les deux structures
réelles de $X$ sont soit $\{E_a, E_{1/a}\}$, soit $\{E'_a, E'_{1/a}\}$, avec
$a > 1$ ; le cas hexagonal est $\{E'_3, E'_{1/3}\}$. Dans le cas carré, elles
sont $\{E_1, E'_1\}$. Le lien avec l'invariant modulaire reste une question
en marge (« $j = \frac12 \log a$ ? ») ; l'invariant est $j(\tau)$, avec le
$\tau$ de la page 98, et il est réel sur les deux familles. La page 95 est un
feuillet de dessins : des sphères découpées par des arcs de grands cercles.

\pagerange{98}{100}
\subsection*{Période, points réels, torseurs (pages 98 à 100)}

Pour $E_a$, une base orthonormée est $(e_0, e_1/\sqrt a)$, d'où la période
$\tau = \sqrt{a}\, i$. Pour $E'_a$, on prend $e_0 = u$ et $e_1 = 2v - u$, avec
$v = \frac12(e_0 + e_1)$, d'où
\[
\tau = \tfrac12\bigl(1 + \sqrt{a}\, i\bigr).
\]

Les points réels sont les $x \bmod \Pi$, $x \in \Pi_{\mathbb{R}}$, tels que
$\sigma x \equiv x \pmod{\Pi}$. Pour $E_a$, $\sigma(x,y) = (x,-y)$, et
$\sigma(x,y) - (x,y) = (0, -2y) \in \Pi$ équivaut à $y \in
\frac12\mathbb{Z}$ : $X(\mathbb{R})$ a \emph{deux} composantes, les images de
$\mathbb{R}e_0$ et de $\frac12 e_1 + \mathbb{R}e_0$.\footnote{Le signe de
$-2y$ est d'une lecture incertaine sur la page ; il est sans effet sur la
conclusion.} Pour $E'_a$, avec $\sigma u = u$ et $\sigma v = u - v$, on a
$\sigma(xu + yv) - (xu + yv) = y(u - 2v)$, qui est dans $\Pi$ si et seulement
si $y \in \mathbb{Z}$ : $X(\mathbb{R})$ est \emph{connexe}. Ainsi $E_a$
correspond, par la section IV, à un anneau (quotient de $X(\mathbb{C})$ par
$\sigma$), et $E'_a$ à un ruban de Möbius.

Reste, dit la page 100 sous une date possible « 9.2. », à classer les
courbes elliptiques réelles inhomogènes : les torseurs $X$ sur
$\operatorname{Spec}\mathbb{R}$ sous les schémas en groupes $E_a$ et $E'_a$.
Soit $\gamma = \mathrm{Gal}(\mathbb{C}/\mathbb{R})$. La suite exacte de
$\gamma$-modules $0 \to \Pi \to V \to E(\mathbb{C}) \to 0$, où $V$, espace
vectoriel réel, a une cohomologie nulle en degrés positifs, donne
\[
H^{1}(\mathbb{R}, E) = H^{1}(\gamma, E(\mathbb{C})) \simeq H^{2}(\gamma,
\Pi) \simeq \Pi^{\gamma}/(1 + \sigma)\Pi .
\]
Pour $E_a$, $\Pi = \Pi_{+} \oplus \Pi_{-}$ et
$H^{2}(\gamma, \Pi) = H^{2}(\gamma, \Pi_{+}) \oplus H^{2}(\gamma, \Pi_{-})
= \mathbb{Z}/2\mathbb{Z} \oplus 0$ : il y a exactement une classe non nulle.
Pour $E'_a$, $\Pi^{\gamma} = \mathbb{Z}u = (1 + \sigma)\Pi$ (car $(1 +
\sigma)v = u$), et $H^{1}(\mathbb{R}, E'_a) = 0$. Les seules courbes
elliptiques réelles vraiment inhomogènes sont donc les torseurs non triviaux
$E_a^{!}$ sous $E_a$, à un paramètre $a \in \mathbb{R}^{*}_{+}$. Un torseur
non trivial n'a pas de point réel. Le quotient de ses points complexes par la
conjugaison est donc une surface conforme non orientable de genre $1$
\emph{sans} bord, une bouteille de Klein : ce sont les surfaces $b)$ de la
page 27.\footnote{La page écrit « sans bord », « sans » paraissant barré
d'un trait, sans qu'aucun mot soit écrit à sa place : elle ne tranche donc
pas. Mais un torseur sous $E_a$ est trivial dès qu'il a un
point réel, et le bord de $X(\mathbb{C})/\sigma$ est l'image de
$X(\mathbb{R})$ : pour un torseur non trivial, il est vide. Les surfaces à
bord de genre $1$ sont les quotients des $E_a$ (anneaux) et des $E'_a$
(rubans de Möbius).}

Ensemble, les sections IV et XII achèvent le tableau des surfaces conformes
de genre $1$ : tores, bouteilles de Klein, anneaux et rubans de Möbius,
chacun avec sa métrique plate canonique.

\end{document}
